\providecommand{\phio}{\phi(0)}
\providecommand{\cR}{\mathcal R}
\providecommand{\Prb}{\mathbb P}
\providecommand{\nm}{\mathrm}
\section{Acquisition and comparison near a two-point law}
\label{j:local:section}
The defect identity has two uses. Contact and a complete Fourier bound
first confine one extremizing marginal to a neighborhood of a two-point
law. We then compare the actual sum with that law's own binomial count.
The second step needs no extremality beyond the simultaneous hypotheses
already obtained in the first. All bounds below include zero excess
$e=C_n-\CE=0$; positive conditional noise and zero noise are treated
separately at the end.

Write $\sigma=\sqrt{10}$, $m=\E X_+=\E X_-$,
$\rho=\E|X|^3$, $\gamma=\E X^3$, $\beta=\E[X|X|]$,
$\mu_j=\E X^j$ and $M_j=\E|X|^j$ (also $\eta_5=M_5$).
In the central gap we use $(-A,B)$ for the endpoints, equivalently
$A=A_-$ and $B=A_+$ in Section~\ref{j:core:section}.
The quantities $J,J_1,J_2,D,L,K_1,K_q$ and $\mathcal R$ are those of
Theorem~\ref{j:core:conservation}; $U=J_1+(\mu_4-3)/6$ below is a
coefficient, distinct from the conditional noise introduced later.
The letter $B$ denotes the positive endpoint until Bernoulli coordinates
are introduced; there its meaning is explicitly reset to a label.


\subsection{Support from the contact inequality}\label{j:local:support:theorem}
\paragraph{The scale correction.}
Write $\phi_0=\phi(0)$.
For $n>1$ put $s_n=\sqrt{n/(n-1)}$ and
$e_n=n(s_n-1)-1/2$.
\begin{lemma}\label{j:local:support:scale}
For every integer $n\ge36$, one has $0<e_n<.384/n$.
\end{lemma}
\begin{proof}
Put $x=1/n$ and $c=.384=48/125$. Direct expansion gives
\[
 (1-x)(1+x/2+cx^2)^2-1
 =x^2\{(2c-3/4)-(c+1/4)x+(c^2-c)x^2-c^2x^3\}.
\]
The bracket is decreasing for $x\ge0$, since $0<c<1$.
At $x=1/36$ its exact rational value is $823/4050000>0$.
Squaring positive quantities gives
$s_n<1+x/2+cx^2$ for $0<x\le1/36$. Similarly
$(1-x)(1+x/2)^2-1=-3x^2/4-x^3/4<0$, which gives the lower bound.
\end{proof}


\paragraph{A signed Gaussian envelope.}
Define, for every real $H,t$ and $n>1$,
\[
 \mathcal H_n(H,t)=n^{3/2}\{\Phi((H-t)/\sqrt{n-1})-
             \Phi(H/\sqrt n)\}
       +nt\phi(H/\sqrt n)+\tfrac H2(t^2-1)\phi(H/\sqrt n).
\]
Write $v=n-1$, $s=\sqrt{n/v}$, $u=|t|$, and
\[
 A=\frac{u^3-3u}{6},\quad B=\frac{u^4-6u^2+3}{24\sqrt v},
 \quad D=\frac{u^5-15u}{120v}.
\]
\begin{lemma}[Fifth-order identity with a uniform finite remainder]
If $u>\sqrt3$, then, with $z=H/\sqrt v$,
\[
 \frac{\mathcal H_n(H,\pm u)}{s^3}
 =\pm A(1-z^2)\phi(z)+B\phi^{(3)}(z)
                  \mp D\phi^{(4)}(z)+R,
\]
where the sign of the remainder need not be fixed and
\[
 |R|\le E^{(5)}_v(u):=
 \frac{51u^6}{14400v^{3/2}}
 +\frac{51(3u^2-2)}{960v^{3/2}}+\frac{u}{8v^2}.
\]
This statement is uniform over every real $H$.
\end{lemma}
\begin{proof}
Let $F_w(x)=\Phi(x/\sqrt w)$, with primes denoting $x$ derivatives.
The heat identity $\partial_w F_w=F_w''/2$ and ordinary Taylor's theorem
give, with unlabelled derivatives at $(v,H)$,
\begin{align*}
 F_{v+1}&=F_v+F_v''/2+F_v^{(4)}/8+R_0,
 &|R_0|&\le\sup_{w\in[v,v+1]}\|F_w^{(6)}\|_\infty/48,\\
 F_{v+1}'&=F_v'+F_v'''/2+F_v^{(5)}/8+R_1,
 &|R_1|&\le\sup_{w\in[v,v+1]}\|F_w^{(7)}\|_\infty/48,\\
 F_{v+1}''&=F_v''+F_v^{(4)}/2+R_2,
 &|R_2|&\le\sup_{w\in[v,v+1]}\|F_w^{(6)}\|_\infty/8.
\end{align*}
Expand $F_v(H-t)$ through degree five, with remainder at most
$u^6\|F_v^{(6)}\|_\infty/720$. In
\[
 \frac{\mathcal H_n}{n^{3/2}}=
 F_v(H-t)-F_{v+1}(H)+tF_{v+1}'(H)
             -\frac{t^2-1}{2}F_{v+1}''(H),
\]
the first three derivative coefficients cancel. The next three are
\[
 \left(\frac t2-\frac{t^3}6\right)F_v'''
 +\frac{t^4-6t^2+3}{24}F_v^{(4)}
 +\left(\frac t8-\frac{t^5}{120}\right)F_v^{(5)}.
\]
The remaining absolute value is bounded by
\[
 \left\{\frac{u^6}{720}+\frac1{48}+\frac{u^2-1}{16}\right\}
       \sup_{w\in[v,v+1]}\|F_w^{(6)}\|_\infty
 +\frac u{48}\sup_{w\in[v,v+1]}\|F_w^{(7)}\|_\infty.
\]
Fourier inversion of the Gaussian, differentiated under its absolutely
integrable Fourier transform, gives
\[
 \|\phi^{(k)}\|_\infty
 \le\frac1\pi\int_0^\infty y^k e^{-y^2/2}\,dy.
\]
Integration by parts yields
$\|\phi^{(5)}\|_\infty\le8/\pi<2.55$ and
$\|\phi^{(6)}\|_\infty\le15\phi_0<6$.
Since $F_w^{(j)}=w^{-j/2}\phi^{(j-1)}$, the displayed remainder
becomes $s^3E^{(5)}_v(u)$ after multiplication by $n^{3/2}$.
The three retained derivative coefficients give the identity.
\end{proof}

\begin{lemma}[A lower tangent bound for the fourth Gaussian derivative]
For every real $z$,
\[
 \phi^{(4)}(z)\ge\phi_0(3-15z^2/2).
\]
\end{lemma}
\begin{proof}
The same Gaussian Fourier integral gives
$\phi^{(4)}(z)=\pi^{-1}\int_0^\infty
y^4e^{-y^2/2}\cos(yz)\,dy$.
Use the globally valid inequality $\cos w\ge1-w^2/2$ inside
this integral. Its fourth and sixth Gaussian moments are $3\phi_0$
and $15\phi_0$, respectively.
\end{proof}

\begin{lemma}[Uniform signed fifth-order envelope]\label{j:local:support:envelope}
Suppose $u>\sqrt3$ and $A,B,D\ge0$. If
$W=A-15D/2>0$, put
\[
 T=A-3D+\frac{9B^2}{4W}.
\]
If $T\ge A/2+2B+3D$, then, for every real $H$ and either sign,
\[
 \mathcal H_n(H,\pm u)\le s^3\{\phi_0 T+E^{(5)}_v(u)\}.
\]
\end{lemma}
\begin{proof}
For positive $t=u$ and $|z|\le1$, divide the retained terms by
$\phi_0$ and use the preceding tangent bound to obtain
\[
 A-3D-(A-15D/2)z^2+3B|z|
 \le A-3D+9B^2/(4W).
\]
Here $(1-z^2)\phi(z)\le(1-z^2)\phi_0$ and
$|\phi^{(3)}(z)|\le3|z|\phi_0$ on this interval.
For $|z|\ge1$ the first term is nonpositive, and the Fourier bounds
$\|\phi^{(3)}\|_\infty<2\phi_0$ and
$\|\phi^{(4)}\|_\infty\le3\phi_0$ give $2B+3D$.
For negative $t=-u$, the first term has global upper bound
$2e^{-3/2}A\phi_0<A\phi_0/2$, and the other terms give
$(2B+3D)\phi_0$. The hypothesis on $T$ covers both exterior cases.
The inequality $2e^{-3/2}<1/2$ follows from four positive terms of
the exponential series at $3/2$. Finally add the uniform remainder.
\end{proof}

\begin{lemma}[A global Gaussian remainder bound]
\label{j:local:ext:global-gaussian-ray}
For every integer $n\ge2$ and all real $H,t$,
\[
 |\mathcal H_n(H,t)|
 <\frac{|t|^3}{15}+\frac{|t|}{5}+\frac14
 <\frac{|t|^3}{15}+\frac{|t|}{5}+\frac12.
\]
\end{lemma}
\begin{proof}
Write $F_w(x)=\Phi(x/\sqrt w)$, $v=n-1$ and $y=H-t$.
The identity defining $\mathcal H_n$ is
\[
 n^{-3/2}\mathcal H_n(H,t)
 =F_v(y)-F_n(H)+tF_n'(H)-\frac{t^2-1}{2}F_n''(H).
\]
Taylor expansion of $F_n(y)$ at $H$ through degree two has remainder
at most $\phi(0)|t|^3/(6n^{3/2})$, since
$\|\phi''\|_\infty=\phi(0)$.
After this expansion the remaining correction is
\[
 \frac12\{F_n''(H)-F_n''(y)\}
 +\frac12\int_v^n\{F_n''(y)-F_w''(y)\}\,dw.
\]
Here we used $\partial_wF_w=F_w''/2$. The first term is at most
$\phi(0)|t|/(2n^{3/2})$ in absolute value. Applying the heat identity
once more bounds the second term by
\[
 \frac14\int_v^n\int_w^n
       |F_r^{(4)}(y)|\,dr\,dw
 \le\frac{\|\phi'''\|_\infty}{8v^2}
 \le\frac1{4\pi v^2}.
\]
The last inequality follows by differentiating the absolutely
convergent Gaussian Fourier integral, which gives
$\|\phi'''\|_\infty\le2/\pi$. Multiply by $n^{3/2}$ and use
$\phi(0)<2/5$. The function $x^{3/2}/(x-1)^2$ decreases for $x>1$,
because its logarithmic derivative is
$3/(2x)-2/(x-1)=-(x+3)/(2x(x-1))<0$. Hence, for $n\ge2$,
\[
 \frac{n^{3/2}}{4\pi(n-1)^2}
 \le\frac{\sqrt2}{2\pi}<\frac14,
\]
using $\sqrt2<3/2$ and $\pi>3$. This proves the stronger bound and
therefore the exact ray majorant used by the original checkers.
\end{proof}

\paragraph{Finite bands and the exact support polynomial.}\label{j:local:support:polynomial}
For integers $36\le N\le L$, use the simultaneous endpoint bounds
\[
 B_+=\frac{u^4-6u^2+3}{24\sqrt{N-1}},\quad
 D_+=\frac{u^5-15u}{120(N-1)},\quad
 D_-=\frac{u^5-15u}{120(L-1)}.
\]
For an infinite index range, set $D_-=0$ instead. Put
\[
 W_*=A-15D_+/2,\qquad
 T_*=A-3D_-+\frac{9B_+^2}{4W_*}.
\]
The proof of Lemma~\ref{j:local:support:envelope} applies uniformly on the whole band
if $W_*>0$ and $T_*\ge A/2+2B_++3D_+$. In particular the favorable
term is frozen at the upper index endpoint, and the adverse terms at
the lower one. No monotonicity of a signed coefficient is assumed.

Let $(P,H)$ be an attained positive closed maximizing pair, standardized
to mean zero and variance one. Write $C=C_n\ge C_E$,
$\rho=\mathbb E|X|^3$, $\beta=\mathbb E[X|X|]$, and
$d=n(C_{n-1}-C_n)_+$. The finite contact/predecessor bound is
\[
 \rho I(t)\le\mathcal H_n(H,t)+
 C\{-|t|^3+3\rho t^2/2+3\beta t+\rho e_n\}+\rho s_nd,
\]
with $I=0$ on the entire closed support. Suppose also, for a finite
band, that $C_n-C_E\ge K/n$, where $K>0$. The finite moment bounds of~\cite[p.~125]{Shevtsova2013}, namely
$C\rho\le .3031(\rho+.646)$ and
$C\rho\le .3322(\rho+.429)$, give
\[
 c=.4095+K/L<C_n,\qquad
 \rho<r=\min\left\{\frac{.3031(.646)}{c-.3031},
                        \frac{.3322(.429)}{c-.3322}\right\}.
\]
For the unrestricted nonnegative case put $K=0$, $c=.4095$.
The sign-moment inequality
$\rho\ge(1+\beta^2)/\sqrt{1-\beta^2}$ supplies $|\beta|<b$ whenever
$(1+b^2)^2>r^2(1-b^2)$ and $0<b<1$. All these statements concern the
same current law. The lower endpoint cannot replace $L$ in $K/L$.

When a preceding resource enclosure supplies sharper actual bounds
$\rho\le r$ and $|\beta|\le b$, these may instead be used in the
same polynomials. Their validity is the only extra premise.

Use rational $\ell^2\le N-1$ and $S^2>N/(N-1)$. For $u=|t|$ define
\[
 Q(u)=-u^3+3ru^2/2+3bu+.384r/N.
\]
First check $Q<0$, before replacing $C$ by its lower bound $c$.
Replace $1/\sqrt{N-1}$ by $1/\ell$ in the adverse remainder and in
the term $2B_+$ of the domination gate. Write $E_*$ for that rational
remainder. Strict negative contact for every $n$ in the band follows
from the polynomial inequality
\[
 \frac25S^3\left\{(A-3D_-)W_*+
              \frac{(u^4-6u^2+3)^2}{256(N-1)}\right\}
 +S^3E_*W_*+\{cQ+rSq\}W_*<0.
\]
The second polynomial gate is $W_*>0$. The third is
\[
 \{A/2-3D_--3D_+-2B_+\}W_*+
         \frac{(u^4-6u^2+3)^2}{256(N-1)}>0.
\]
This is the cleared-denominator version of the domination hypothesis.
The bound on $T_*$ is positive, so the replacement $\phi_0<.4$ and
$s_n<S$ both have their correct direction.

The exact kernel tests these polynomials on all adjacent closed intervals
of length $.05$ from $R$ to $4.5$. Positive Bernstein coefficients prove
their signs on every point of each cell, including endpoints. On the
rays $u\ge4.5$ use Lemma~\ref{j:local:ext:global-gaussian-ray}, which gives the same bound
$|\mathcal H_n(H,t)|\le u^3/15+u/5+1/2$.
The resulting cubic and its first three derivatives are strictly negative
at $4.5$, as is the coefficient $Q$ before replacing $C$ by $c$.
This excludes the whole rays. Thus, if $d\le q$, actual support is in
$(-R,R)$, including exclusion of arbitrarily small exterior atoms and
exterior accumulation points.


\paragraph{The support rows used in this proof.}
For the all-index local argument, take $K=0$, $R=3.84$, $q=.19$ on
\[
[45,49],\ [50,59],\ [60,89],\ [90,134],\ [135,199],\ [200,\infty).
\]
Only indices at least 47 are used. The last row uses $D_-=0$;
there is no upper index restriction. The exact checker reconstructs
all closed $u$ cells and the two exterior rays. Missing an index band
or the infinite row is an acceptance failure.

\paragraph{One bounded-threshold series.}\label{j:local:support:series}
Put $c_k=(-1)^k/[2^kk!(2k+1)]$, and let $(x)_j$ denote the rising
factorial. Define the following rational polynomials:
\begin{align*}
 P_j(H,t)={}&\sum_{k=0}^{j+1}
 \frac{c_k(k+1/2)_{j+1-k}}{(j+1-k)!}(H-t)^{2k+1}
       -c_{j+1}H^{2j+3}\\
 &+\frac{(-1)^{j+1}tH^{2j+2}}{2^{j+1}(j+1)!}
       +\frac{(-1)^jH^{2j+1}(t^2-1)}{2^{j+1}j!}.
\end{align*}
In particular,
\[
 P_0=(t^3-3t)/6,\qquad
 P_1=\frac{3t-t^3}{4}H^2+\frac{t^4-6t^2+3}{8}H
                  +\frac{-t^5+10t^3-15t}{40}.
\]
Fix $h>0$, $E(x)=(193/71)^{\lceil x\rceil}$, and
\[
 B_r=\frac{7(h+r)}6E((h+r)^2/6)
       +\left(h+r+\frac{h(r^2+1)}8\right)E(h^2/8).
\]
\begin{lemma}\label{j:local:support:finite-series}
For every $n\ge4$, $|H|\le h$ and $|t|\le r$,
\[
 \left|\mathcal H_n/\phi(0)-\sum_{j=0}^5P_j/n^j\right|
       \le4^7B_r/n^6.
\]
\end{lemma}
\begin{proof}
Put $z=1/n$. Integration of the Gaussian exponential series gives
$\mathcal H_n/\phi(0)=S(z)/z$, where
\begin{align*}
 S(z)={}&\sum_{k\ge0}c_k\{(H-t)^{2k+1}(1-z)^{-k-1/2}
                                  -H^{2k+1}\}z^k\\
 &+te^{-H^2z/2}+\tfrac12zH(t^2-1)e^{-H^2z/2}.
\end{align*}
Its constant coefficient vanishes, and expansion of the binomial factor
gives $[z^{j+1}]S=P_j$. At $z_0=1/4$, the sum obtained by replacing
all coefficients by their absolute-value majorants is at most $B_r$:
use $(1-z_0)^{-1/2}<7/6$, $1/(2k+1)\le1$, and $e<193/71$.
The latter bound follows, for example, by summing the exponential series
through degree six and bounding the rest by a geometric series with
ratio $1/8$. All products and rearrangements are therefore absolutely
convergent. The discarded terms have degrees at least seven in $S$,
so their sum at $0\le z\le z_0$ is at most $(z/z_0)^7B_r$.
Divide by $z$ to get the result. This proves a finite uniform bound,
including every threshold in the stated interval and every $n\ge4$.
\end{proof}

\paragraph{A common contact test for radius and strip refinement.}
\label{j:local:support:panels}
Let $N\ge36$, $E,q\ge0$ and $C_0>0$. Suppose $n\ge N$,
$0\le C-C_E\le E/n$, $\delta_n:=n(C_{n-1}-C_n)_+\le q$, $C\le C_0$ and
$|X|\le R_0$. A previously acquired row gives
\[
 \rho\le r,\quad 0\le b_l\le\beta\le b_u<1,\quad |H|\le h.
\]
Define $Q(t)=-|t|^3+3\rho t^2/2+3\beta t$ and
$L_P(t)=C_EQ(t)+\phi(0)(t^3-3t)/6$.
For $t>0$, upper polynomials for $L_P(t),Q(t)$ are
\[
 (1/15-.4095)t^3+1.5(.40974)rt^2+(3(.40974)b_u-.3988/2)t,
 \quad -t^3+1.5rt^2+3b_ut.
\]
For $t<0$, write $u=|t|$ and use
\[
 (-.3988/6-.4095)u^3+1.5(.40974)ru^2+(.2-3(.4095)b_l)u,
 \quad -u^3+1.5ru^2-3b_lu.
\]
These coefficient choices are valid because $b_l\ge0$ is checked first.
On any closed cell of one sign, let $L_J$ be an upper bound for the
first polynomial and let $Q_J\ge0$ bound the positive part of the second.
For $j=1,2,3$, take $m_j$ to be the positive part of the largest tensor
Bernstein coefficient of $P_j$ on $[-h,h]\times J$. For $j=4,5$, the
absolute monomial bound
$\sum_{a,b}|[H^at^b]P_j|h^aR_J^b$ suffices, where
$R_J=\max_{t\in J}|t|$. The support contact is then bounded above by
\begin{equation}\label{j:local:support:paneltest}
 L_J+\frac{E Q_J}{N}
 +.4\left(\sum_{j=1}^5\frac{m_j}{N^j}
                +\frac{4^7B_{R_J}}{N^6}\right)
 +r\left(qS+\frac{.384C_0}{N}\right),
 \qquad S>\sqrt{N/(N-1)}.
\end{equation}
Indeed substitute the finite series in contact, split $C=C_E+(C-C_E)$,
and use the scale estimate $e_n<.384/n$. Every error retained in
\eqref{j:local:support:paneltest} is nonnegative and decreases with $n$.
A strictly negative rational result therefore excludes the entire cell
for every integer $n\ge N$, including atoms and accumulation points.
The uniform-tail specialization is $(N,E,q,C_0)=(47,.7,.19,.425)$;
its score cap follows from $C<.40974+.7/47<.425$.
The finite bands use their own acquired $E,q,C_0$ at the corresponding
floor, with exactly the same formula.

For a coordinate strip, conditional centering gives
\[
 X=\ell(Y-p),\quad
 \ell=\sqrt{(1-\theta)/(p(1-p))},\quad
 \beta=(1-\theta)(1-2p)+\kappa,\quad |\kappa|\le\theta.
\]
If $p\in[p_l,p_u]\subset(0,1/2)$ and $\theta\le\Theta$, put
$\ell_l\le\sqrt{(1-\Theta)/(p_u(1-p_u))}$.
The forbidden strip $Y\notin(-.55,1.55)$ maps into
\[
 X\le-\ell_l(p_l+.55),\qquad X\ge\ell_l(1.55-p_u).
\]
Apply \eqref{j:local:support:paneltest} between these thresholds and $\pm R_0$.
An interval beyond the known radius is empty; an equal endpoint is a
singleton that is still checked. If only an upper $\beta$ bound is given,
use $b_l=(1-\Theta)(1-2p_u)-\Theta$, after checking its nonnegativity.

The selected consequences used in the resource schedule are:
\[
\begin{array}{l|c|c|r}
\text{input rows}&\text{prior radius}&\text{conclusion}&\text{closed cells}\\\hline
10\text{ central preparation rows}&3.84&|X|<2.5&1080\\
45\text{ central rows}&2.5&-.55<Y<1.55&1431\\
38\text{ left rows}&2.7&-.55<Y<1.55&3073
\end{array}
\]
The respective largest contact bounds are below $-.084$, $-.097$ and
$-.01658$. These implications acquire no row by themselves: each is
used only after its input row is proved by the preceding resource step.
The left radii are applications of the all-threshold support polynomial,
with $(N,L,K,q)=(47,\infty,0,.19)$ and the actual moment caps
\[
 (r,b,R)=(1.390,.430,3),\qquad (1.2563,.3926,2.7).
\]
Thus the same heat proof supplies the initial radius, both left
contractions, and the five finite-base rows. The bounded-threshold
series supplies the central contraction and both coordinate strips.




\subsection{Acquiring simultaneous resources for one marginal}\label{j:local:sign}
For a ratio maximizer $C=C_n\ge\CE$, put $e=C-\CE$ and $w=H^2/(2n)$.
The core identity is
\[
 ne\rho+n\CE\rho(1-e^{-w})+n\phi(0)e^{-w}D/6+n^{3/2}L=\mathcal R.
\]
The first complete raw enclosure must give $\mathcal R<r<.4N$, for the
actual $n\ge N$. Since $C\rho\ge\CE>.4$, this licenses
Corollary~\ref{j:core:license}, including
\[
 ne\rho+n\phi(0)D/6+n^{3/2}L\le\mathcal R,\quad
 e^{-w}\ge1-5r/(2n),\quad nD\le15000r/997.
\]
The proof uses $e\ge0$ and never divides by $e$.
If a previous pass gives $nD\le Nd$, subtracting the favorable phase
term gives the sharper shared budget
\begin{equation}\label{j:local:budget}
 ne\rho+\kappa nD+n^{3/2}L\le\mathcal B(H),\qquad\kappa=997/15000,
\end{equation}
by increasing each raw quadratic curvature by
\begin{equation}\label{j:local:re:extra}
 q_{\rm extra}(\rho)=\tfrac12(1-5r/(2N))(.4095\rho-d/15)>0.
\end{equation}
Indeed $1-e^{-w}\ge we^{-w}$, $D\le d$, $\phi(0)<.4$ and
$\CE>.4095$. The positivity gate is checked before this use.
The field called \texttt{Qbudget} already contains this addition.
For a cell with $D\ge d_l\ge0$ and $\rho\ge\rho_l>0$, the same identity
also gives the score update
\[
 C\le .40974+\{r/N-.3988(1-5r/(2N))d_l/6\}/\rho_l.
\]
A negative numerator excludes the cell by $e\ge0$. Otherwise the numerator
increases with $1/n$, proving this bound at every $n\ge N$.


\paragraph{One positive-measure rule.}
Suppose zero lies in the gap $(-A,B)$, $A,B>0$. Let
\[\begin{gathered}
 p=\Pp(X>0),\quad q=1-p,\quad m=\E X_+=\E X_-,\quad V=\E[X\mid\mathbf1_{X>0}],\\
 W=X-V,\quad\theta_\pm=\E[W^2;\pm X>0],\quad\theta=\theta_++\theta_-.
\end{gathered}\]
Then $m^2=(1-\theta)pq$, and the conditional means have magnitudes
$\bar A=m/q$, $\bar B=m/p$. Set
$\Delta_A=\bar A-A\ge0$, $\Delta_B=\bar B-B\ge0$.
We also use the unnormalised endpoint displacements
\[
 M_+=\E[(X-B)\mathbf1_{X>0}]=p\Delta_B,\qquad
 M_-=\E[(-X-A)\mathbf1_{X<0}]=q\Delta_A.
\]

\begin{lemma}[Positive-measure pricing]\label{j:local:re:pricing}
For $x=|X|$, $a=|V|$, and a function $h$ constant on each sign class,
\begin{equation}\label{j:local:re:conditional}
 \E[h(|X|^j-|V|^j)]
 =\E\left[hW^2\sum_{k=0}^{j-2}(k+1)a^kx^{j-2-k}\right],\quad j\ge2.
\end{equation}
Also, for $d\mu_0=(|X|-1)^2(|X|+1/2)\,dP$,
\begin{align*}
 \mu_0(\R)&=\rho-1,\\
 M_4-1&=\int\frac{(x+1)^2}{x+1/2}\,d\mu_0,\\
 M_5-1&=\int\frac{x^3+2x^2+3x+3/2}{x+1/2}\,d\mu_0,\\
 K_1+3\rho\beta&=\int
  \frac{\operatorname{sgn}(X)x(x+2-3\rho/2)}{x+1/2}\,d\mu_0.
\end{align*}
For any nonnegative resources $z_i$ satisfying $\sum a_i z_i\le r$,
$a_i>0$, a signed linear cost obeys
\begin{equation}\label{j:local:re:simplex}
 \sum b_i z_i\le r\max\{0,b_1/a_1,\ldots,b_k/a_k\}.
\end{equation}
\end{lemma}
\begin{proof}
Divide $x^j-a^j-ja^{j-1}(x-a)$ by $(x-a)^2$; the linear term has
zero conditional mean. The radial identities follow by multiplying
out and using $\E X=0$, $\E X^2=1$. Finally each $b_i$ is at most
$a_i$ times the displayed maximum, and all $z_i$ are nonnegative.
\end{proof}

This one rule accounts for the moment, signed Taylor, endpoint movement,
outer-gap, and confinement-error estimates. In particular, no two costs
may each spend the full residual and then be added.



The gap identity and simultaneous defect budget are
\eqref{j:core:gap-defect} and \eqref{j:core:simultaneous}, with
$(A_-,A_+)=(A,B)$ and $(\Delta_-,\Delta_+)=(\Delta_A,\Delta_B)$.
The outer-gap estimate \eqref{j:core:outer-gaps} bounds the same
$T_\pm,R_\pm,J_o$. One resource pays for all their charges together.

\paragraph{A state, its update, and its exported resource.}
A state for a closed probability interval $I=[l,u]$ records
\[
 |X|\le R,\quad C\in[c_0,c_1],\quad\mathcal R\le r,
 \quad nD\le Nd,\quad\theta\le\Theta<1,
 \quad A\ge A_l>0,\quad B\ge B_l>0,
\]
together with simultaneous signed-moment and coefficient intervals.
The distinction between $nD\le Nd$ and $D\le d$ is essential.
Let $s_+=\max_Ip(1-p)$, $s_-=\min_Ip(1-p)$ and
\begin{gather*}
 b_I=\frac{(6+2\sigma)(1-\Theta)^{3/2}}{\sqrt{s_+}}
                 \operatorname{dist}(p_*,I)^2,
 \qquad r_I=d-b_I,\\
 K_+=3(\sigma+1)B_l,\quad K_-=3(\sigma+3)A_l,
 \quad K=\min(K_+,K_-),\quad m_l=\sqrt{(1-\Theta)s_-}.
\end{gather*}
If $r_I<0$, exclude the entire closed cell. Otherwise \eqref{j:core:simultaneous}
gives, simultaneously,
\begin{equation}\label{j:local:re:residual}
 6m(A_l+B_l)(\Delta_A+\Delta_B)+K_+\theta_++K_-\theta_-\le r_I.
\end{equation}
Thus the actual law satisfies the intersected bounds
\begin{gather}
 \theta\le\min\{\Theta,r_I/K\},\notag\\
 A\ge\max\left\{A_l,\sqrt{(1-\Theta)l/(1-l)}-
                  \frac{r_I}{6m_l(A_l+B_l)}\right\},\notag\\
 B\ge\max\left\{B_l,\sqrt{(1-\Theta)(1-u)/u}-
                  \frac{r_I}{6m_l(A_l+B_l)}\right\}.\label{j:local:re:geometry}
\end{gather}
Each right side is evaluated from established bounds; sequential
improvements may then be used by the next expression. A nonpositive
endpoint estimate fails this update, rather than excluding a law.
Initialization uses the same formulas with the valid lower sum
$A+B\ge\sqrt{2-d}$.

In Bernoulli coordinates set $Y=B_0+U$, $\E[U\mid B_0]=0$,
$B_0\sim\mathrm{Bernoulli}(p)$, $X=\ell(Y-p)$ and
$d_Y=(\Delta_A+\Delta_B)/(\bar A+\bar B)$.
Then $\ell^2=(1-\theta)/(pq)$ and $\lambda=n\E U^2=npq\theta/(1-\theta)$.
The \emph{weighted} versions of \eqref{j:local:re:residual} prove
\begin{equation}\label{j:local:re:exports}
 n\theta\le\frac{Nr_I}{K},\quad
 nd_Y\le\frac{Nr_I}{6(1-\Theta)(A_l+B_l)},\quad
 \lambda\le\frac{Ns_+r_I}{(1-\Theta)K}.
\end{equation}
Indeed $nD\le Nd$ whereas $nb_I\ge Nb_I$.
These exports are derived from the weighted resource; multiplying an
unweighted variance ceiling by $N$ would not prove them.

For $j=3,4,5$ set
\[
 L_3(v)=R+2v,\quad L_4(v)=R^2+2Rv+3v^2,
 \quad L_5(v)=R^3+2R^2v+3Rv^2+4v^3.
\]
Put $v_+=\sqrt{(1-l)/l}$, $v_-=\sqrt{u/(1-u)}$ and
$k_j=[1-(1-\Theta)^{j/2}]/\Theta$, with $k_j=j/2$ at $\Theta=0$.
For $u<1/2$, valid Bernoulli moment ceilings are, with $s=l(1-l)$,
\[
 B_3^*=(1-2s)/\sqrt s,\quad B_4^*=1/s-3,
 \quad B_5^*=(1-4s+2s^2)/s^{3/2}.
\]
They decrease towards $p=1/2$. Convexity places $(1-\theta)^{j/2}$
below its endpoint chord on $[0,\Theta]$; positive-measure pricing gives
\begin{equation}\label{j:local:re:moments}
 M_j\le B_j^*+r_I\max\left\{0,
 \frac{L_j(v_+)-k_jB_j^*}{K_+},
 \frac{L_j(v_-)-k_jB_j^*}{K_-}\right\}.
\end{equation}
The same rule, using
$T_+\le2mv_+\Delta_B+L_3(v_+)\theta_+$ and its reflection, bounds
\begin{equation}\label{j:local:re:outerresource}
 W_o:=(\sigma-1)T_++(\sigma+1)T_-
 \le r_I\max\left\{\begin{gathered}
 \frac{(\sigma-1)v_+}{3(A_l+B_l)},
 \frac{(\sigma+1)v_-}{3(A_l+B_l)},\\
 \frac{(\sigma-1)L_3(v_+)}{K_+},
 \frac{(\sigma+1)L_3(v_-)}{K_-}\end{gathered}\right\}.
\end{equation}
Intersect every new bound with the previous one. The versions which
merge $\theta_+$ and $\theta_-$ are obtained by replacing both costs by
their maximum and both prices by their minimum; they need no new proof.

\paragraph{The enclosure theorem.}\label{j:local:re:operator}
On each cell, the formula atlas bounds the four coefficients in
\eqref{j:core:remainder}. The result, for each actual sign $s=\operatorname{sgn}H$,
is a finite family of quadratics
\begin{equation}\label{j:local:re:quadratics}
 \mathcal R\le\max_j(A_j+F_{j,s}|H|-Q_j^{\rm raw}H^2),
 \qquad Q_j^{\rm raw}>0.
\end{equation}
The maxima retain all score--moment corners. In the initial signed
cover, intervals contain both signs of $\beta,\gamma$; no orientation is
discarded. In the probability cover, the central integral bound is a
maximum of affine functions of $\rho$, hence convex. The other
$\rho$-dependence is affine or the maximum of convex quadratics. For
fixed $H$, the score dependence is convex. Endpoints are therefore taken
\emph{before} completing the square.

\begin{theorem}[Finite resource enclosure]\label{j:local:acquisition}
Consider a finite ordered schedule of the following operations:
\begin{enumerate}
\item Split a cell into an explicitly covered family of closed cells.
\item Intersect its enclosures with an inequality proved in the atlas.
\item Exclude a cell by an empty necessary interval or a strictly
negative nonnegative resource.
\item Apply a separately proved support or confinement implication,
checking all its hypotheses against this cell's already acquired data.
\item Export a raw, weighted, or signed phase bound by
\eqref{j:local:re:quadratics}, \eqref{j:local:budget}, and \eqref{j:local:re:exports}.
\end{enumerate}
If exact outward arithmetic checks every domain, corner and cover,
every output contains all attaining laws satisfying the input. The
conclusion holds for all integers $n\ge N$ wherever the formula's
monotonicity gates are checked. It includes zero conditional variance,
singleton intersections and arbitrarily small atoms.
\end{theorem}
\begin{proof}
Induct over the schedule. A split preserves the union, intersections use
necessary inequalities, and an excluded cell contains no feasible law.
Each callback consumes only earlier hypotheses for that same law.
For a scalar raw output, use
$fh-qh^2\le(f_+)^2/(4q)$ for $h\ge0$, $q>0$; using $f^2$ is a valid
weaker bound. For a resource output first add \eqref{j:local:re:extra} to each
curvature. A negative maximum excludes the cell; a nonnegative one may
be divided by $\kappa N$ to give its next defect parameter.
The phase bound follows from the positive root of
$A+fh-qh^2\ge0$, replacing $A$ by $A_+$ if necessary.
A strict output is exported only when the certificate retains a strict
upper gap, such as its strict Fourier or target enclosure; closed-domain
coverage does not change strictness. All constants are frozen before evaluating
corners. Thus no iteration
step relies on an unproved fixed point. The stated index uniformity is
that of the displayed finite formulas, not a fixed-law limit.
\end{proof}

For the Fourier cancellation estimate only, reflect $(X,H)$ if needed:
$I_0(-X,-H)=I_0(X,H)$. Both odd coefficients change sign, and
$\int e^{-x^2/2}|zx^2-cx^4|\,dx$ is unchanged by $(z,c)\mapsto(-z,-c)$.
Consequently a lower endpoint used there bounds $|\gamma|$, not necessarily
the original signed $\gamma$. The actual contact, gap and phase data keep
their original orientation. A useful additional signed lower bound is
\begin{equation}\label{j:local:re:gamma}
 \gamma\ge(1-\Theta)^{3/2}\frac{1-2u}{\sqrt{u(1-u)}}
       -\left(R+2\sqrt{\frac{u}{1-u}}\right)\Theta.
\end{equation}
To prove it expand $\E X^3=\E V^3+\E[W^2(X+2V)]$; the positive
residual is nonnegative and the negative residual has the displayed
lower bound. The binary term decreases in both $p<1/2$ and $\theta$.
Thus $\max(0,\text{any signed lower bound})$ is a valid lower bound for
$|\gamma|$ inside this auxiliary Fourier estimate.



\paragraph{Signed moments with a common first moment.}\label{j:local:atlas}\label{j:local:moments}
Put $a_\pm=(1\pm\beta)/2$, $b_\pm=(\rho\pm\gamma)/2$.
Cauchy gives $m^2\le a_+a_-$ and $b_\pm m\ge a_\pm^2$.
Consequently
\[
 \frac{(1+\beta)^{3/2}}{\sqrt{1-\beta}}-\rho
 \le\gamma\le
 \rho-\frac{(1-\beta)^{3/2}}{\sqrt{1+\beta}},\qquad
 m\ge\max\left\{\frac{1+\beta^2}{2\rho},
                  \frac{a_+^2}{b_+},\frac{a_-^2}{b_-}\right\}.
\]
These statements also hold with a zero atom. Projection of $X|X|$ onto
the span of $X$ and $|X|$ gives
\[
 \mu_4\ge\rho^2+\frac{(\gamma-\rho\beta)^2}{1-\beta^2}.
\]
For completeness, this follows by minimizing the nonnegative expectation
$E(X|X|-uX-v|X|)^2$ over $u,v$; its Gram matrix has diagonal entries
one and off-diagonal entry $\beta$, and is positive definite.

The signwise moment matrices below, valid also for different radii,
provide the bounded-support fourth- and fifth-moment contractions.
A second useful upper bound for $m$ follows by integrating
$(x-a)^2(R-x)\ge0$ for $x=|X|\in[0,R]$. For every fixed $a>0$,
\[
 2m\le\frac{R-\rho+2a+Ra^2}{a(2R+a)}.
\]
The cover uses only the eleven rational choices $a=2/10,\ldots,12/10$.

Write $\delta=\gamma-\rho\beta$. Cauchy and a square completion on
each sign give the signed constant bound
\begin{align*}
 \frac\rho2\mu_4-\frac14\eta_5
 &\le\frac{\rho^3}{4}
 -\frac{(\rho+\gamma)(\delta_+)^2}{8(1+\beta)^2}
 -\frac{(\rho-\gamma)((-\delta)_+)^2}{8(1-\beta)^2}.
\end{align*}
Indeed on either sign, write $a=\E X_s^2$, $b=\E X_s^3$,
$c=\E X_s^4$, $f=\E X_s^5$. Cauchy gives $c\ge b^2/a$,
$f\ge c^2/b$, and
$\rho c/2-f/4\le\rho^2b/4-(c-\rho b)^2/(4b)$.
Distance to the half-line $c\ge b^2/a$ proves the stated penalty
on summing both signs. Independently,
$\mu_4^2\le\rho\eta_5$ and the projection bound above give
\[
 \frac\rho2\mu_4-\frac14\eta_5
 \le\frac{\rho^3}{4}
 -\frac{(\gamma-\rho\beta)^4}{4\rho(1-\beta^2)^2}.
\]
The larger of the two nonnegative deductions is retained.

\paragraph{An actual central gap without assuming \texorpdfstring{$D<2$}{D below 2}.}
If zero belongs to the support, the signwise gap calculation gives $J\le\rho$.
Thus $D\ge\sqrt{10}\rho-\gamma\ge(\sqrt{10}-1)\rho$.
Whenever $D<(\sqrt{10}-1)\rho$, there is therefore a central gap
$(-A,B)$, $A,B>0$, with no zero atom. Put $p=P(X>0)$, $q=1-p$,
$T_+=b_+-mB^2$ and $T_-=b_--mA^2$. The exact gap decomposition is
\[
 D=mQ+(\sigma-1)T_++(\sigma+1)T_-+3(T-J_o),\qquad
 Q=(\sqrt{\sigma+1}A-\sqrt{\sigma-1}B)^2,
 \quad\sigma=\sqrt{10}.
\]
All terms are nonnegative. Hence $T\le D/(\sigma-1)$ and
\[
 (A+B)^2\ge A^2+B^2\ge\frac{\rho-D/(\sigma-1)}m.
\]
Moreover $A\le a_-/m$, $B\le a_+/m$ and $J\le\rho+2m$.
The last bound gives the compulsory defect
\[
 D\ge \sigma\rho-\gamma-6m.
\]
The inequalities $T-J_o\ge2BR_++2AR_-$ give the useful additional
endpoint constraint
\[
 D\ge(\sigma-1)b_++6a_+B-(\sigma+5)mB^2.
\]
If $D<(\sigma-1)b_+$, this downward quadratic forces $B$ to be at
least its positive root. The reflected constraint has $a_-,b_-,A$
and the coefficients $\sigma+1,\sigma+7$. In a closed moment box,
positive coefficients are rounded down, negative coefficients up, and
the lower square-root endpoint is used in the resulting root bound.

The actual opposite-sign event has probability $2pq$.
The first two moment constraints give lower bounds for both $p$ and $q$.
For independent copies one has exactly
\[
 E[(X-X')^2\mid XX'<0]
 =\frac{1+\beta(1-2p)+4m^2}{2pq}
 =\frac{1+4m^2+\beta}{2p}
  +\frac{1+4m^2-\beta}{2q}.
\]
For an upper bound, use the upper $m$ endpoint and both endpoints of
$p$ and $\beta$. The expression is convex in $p$ and affine in $\beta$.
This improves the former $2/P(XX'<0)$ bound inside the complete HIGH
estimate. Both copies still have the same actual marginal.

\paragraph{Preserving central and outer integral costs.}
Let
\[
 P_1(A,B)=B^3/3+AB^2/2-A^3/6,\quad
 P_2(A,B)=B^4/4+AB^3/3-A^4/12,\quad G=\rho P_1-P_2.
\]
Then $J_{1,c}=mP_1$, $J_{2,c}=mP_2$. Since
$\partial_BG=B(A+B)(\rho-B)$,
\[
 G(A,B;\rho)\le G(A,\rho;\rho)
 =\rho^4/12+\rho^3A/6-\rho A^3/6+A^4/12.
\]
The last expression increases with $A$ because its derivative is
$(\rho-A)^2(\rho+2A)/6$, and is convex in $\rho$ because its second
$\rho$ derivative is $\rho^2+\rho A$. Its closed-rectangle maximum is
therefore bounded using the upper $A$ endpoint and both $\rho$ endpoints.

Positive outer gaps satisfy
$(t-a)(\rho-t)\le(\rho-B)_+^2/4$. On negative outer gaps the positive
part is bounded by $(R-\rho)^2/4$ after extracting $|t|$.
Integration gives
\[
 3(\rho J_1-J_2)_o
 \le\frac38\{(\rho-B)_+^2T_++(R-\rho)^2T_-\}.
\]
The available weighted outer budget is $D-mQ$, not the entire $D$.
The cover intersects this budget with the individual bounds
$T_+\le b_{+,u}-m_lB_l^2$ and $T_-\le b_{-,u}-m_lA_l^2$.

For the signed phase term $U=J_1+(\mu_4-3)/6$, retain
\[
 -RT_-/2\le J_{1,o}\le RT_+/2.
\]
The central minimum of $P_1$ occurs at $B_l$ and an $A$ endpoint;
its maximum occurs at $B_u$ and $A$ clipped to the interval at $A=B_u$.
A nonnegative polynomial bound is multiplied by $m_u$ for an upper
bound and $m_l$ for a lower bound; a negative one uses the opposite
endpoint. For example, the binary endpoints $A=3/2$, $B=2/3$ make this
maximum negative; the sign switch is necessary even when $D<2$.
The fourth-moment sign intervals also sharpen
$K_1=E[X|X|^3]-3\beta-3\rho\gamma/2$.
The complete $K_1,U,\phi$ rectangles are used in the signed expression
$|CK_1+\phi U|$, without replacing it by two separate absolute values.


\paragraph{A complete coarse bound before a central gap is licensed.}
Always $-\mu_4/6\le J_1\le\mu_4/2$ and
$-3J_2-\eta_5/4\le0$. On a gap of length $h=b-a$, these follow from
\[
 \int_a^b(t-a)t\,dt+\tfrac12\int_a^bt^2dt=hb^2/2,\quad
 \tfrac32\int_a^bt^2dt-\int_a^b(t-a)t\,dt=h(2a+b)^2/6,
\]
and, on a negative gap $(-v,-u)$, from
$\tfrac13\int_u^v x^3dx-\int_u^v(v-x)x^2dx=u^3(v-u)/3$.
For $1\le r\le\rho$ and actual support radius $R\le6$, set
\[
 \Gamma(R,r)\ge\max_{r\le v\le R}
 \left\{-\frac{v^2}4+\frac{(R+r)v}3-\frac{Rr}2+
                      \frac{r^3(2R-r)}{12v^2}\right\}_+.
\]
Then
\[
 3(\rho J_1-J_2)\le\frac32\Gamma(R,r)(\rho-\gamma)
                              +\frac38\rho^2J.
\]
On negative gaps extract $x=|t|$ and discard the nonpositive part
$x<\rho$. Replacing $\rho$ by $r$ and the outer endpoint by $R$
gives $\int_r^v(R-x)x(x-r)\,dx$. Dividing this integral by $v^2$
gives exactly the displayed expression. Signwise Tonelli now charges
$\Gamma\E X_-^3=\Gamma(\rho-\gamma)/2$.
On positive gaps $(t-a)(\rho-t)\le\rho^2/4$, giving $\rho^2J/8$.
The implementation uses 64 closed $v$-intervals: on each, maximize the
concave quadratic at its clipped vertex $2(R+r)/3$ and the positive
reciprocal term at the left endpoint. This proves its interval bound.
The alternative bound $\Gamma\le5/2$ follows at $r=1,R=6$ by
$3v^4-28v^3+66v^2-11>0$ on $[1,6]$: the endpoint values are $30,205$,
and the interior critical values are $(85\mp35\sqrt5)/2>0$.
These monotone enclosing parameters cover all $r\ge1,R\le6$.
\paragraph{Moment and gap coefficients for a whole cell.}
\label{j:local:mom:conditional-boxes}

The following coefficient rules are collected here to make the finite
calculation transparent. Every rule is an inequality for an actual law;
there is no extremal-law assumption inside a cell.
Write $v=1-\theta$. For $p\in[l,u]\subset(0,1)$ set
\begin{gather*}
 s_0=\min\{l(1-l),u(1-u)\},\quad
 v_* =\max\{\sqrt{u/(1-u)},\sqrt{(1-l)/l}\},\\
 r_B=(1-2s_0)/\sqrt{s_0},\quad
 g_B=\max_{t=l,u}|1-2t|/\sqrt{t(1-t)},\quad b_B=\max_{t=l,u}|1-2t|,\\
 L_3=R+2v_*,\quad L_4=R^2+2Rv_*+3v_*^2,\\
 L_5=R^3+2R^2v_*+3Rv_*^2+4v_*^3,\quad
 e_B=r_B\frac{1-4s_0+2s_0^2}{s_0(1-2s_0)}.
\end{gather*}
The polynomials $T_3=x+2a$, $T_4=x^2+2ax+3a^2$ and
$T_5=x^3+2ax^2+3a^2x+4a^3$ are nonnegative. Their maxima at
$x=R$, $a=v_*$, in Lemma~\ref{j:local:re:pricing}, give, if $\theta\le\Theta<1$,
\begin{align}
 \rho&\le r_B+(L_3-r_B)\Theta,&
 |\gamma|&\le g_B+(L_3-g_B)\Theta,\notag\\
 |\beta|&\le b_B+(1-b_B)\Theta,&
 \mu_4&\le(1-\Theta)^2(1/s_0-3)+L_4\Theta,\notag\\
 \eta_5&\le(1-\Theta)^2e_B+L_5\Theta.&&\label{j:local:mom:boxes}
\end{align}
Replacing $\theta$ by $\Theta$ uses the explicitly checked inequalities
$L_3>r_B,g_B$, $L_4>2(1/s_0-3)$ and $L_5>2e_B$.
The exact binary factors are $v^{3/2},v,v^2,v^{5/2}$; the upper bounds
use $v^{3/2}\le v$ and $v^{5/2}\le v^2$. Thus cells crossing $1/2$ and
zero conditional variances are included. The same cubic expansion, in the
later domain $p\in[.37,.47]$, $A\ge.75$, $B\ge1.05$, gives
\begin{equation}\label{j:local:mom:rho-lower}
 \rho\ge v^{3/2}\rho_B(p)+3\min(A,B)\theta
       \ge\rho_B(p)+.45\theta,\qquad
 \rho_B(p)=\frac{1-2pq}{\sqrt{pq}}.
\end{equation}
Indeed $\rho_B<1.11$ there and $v^{3/2}\ge1-3\theta/2$.
Also $\beta\ge1-2p-2(1-p)\theta$.

For the signed coefficient
$K_1=\E[X|X|^3]-3\beta-3\rho\gamma/2$, define
\[
 K(p,v)=(1-2p)\{(3v^3/2-v^2)/(pq)-3v^3+v^2+3v\}.
\]
Direct substitution gives $K_1(V)=-K(p,v)$, with the coefficient $3$
unchanged. For $v\ge2/3$, $K$ decreases in $p$: both coefficients in
braces are nonnegative and $(1-2p)/(pq)$ decreases. Its derivative in $v$
has the sign of $1-2p$, since the derivative of the braces is at least
$9v^2-6v+3>0$. Hence both variance endpoints are retained. The error is
\begin{equation}\label{j:local:mom:K-error}
 |K_1(X)-K_1(V)|\le
 \Theta\{L_4+3+\tfrac32L_3(R_\rho+g_B)\}
\end{equation}
when $\rho\le R_\rho$. This follows from the conditional expansions and
$\rho\gamma-\rho_V\gamma_V=
\rho(\gamma-\gamma_V)+\gamma_V(\rho-\rho_V)$.

Here are the remaining signed charges. They also explain the two
alternative bounds retained in the initial, coarse moment cover:
\begin{align}
 -\mu_4/6&\le J_1\le\mu_4/2,&
 -3J_2-\eta_5/4&\le0,                                  \label{j:local:mom:gap-signs}\\
 \rho J_1-J_2&\le\tfrac54(\rho-\gamma)+\tfrac18\rho^2(\rho+1),
                                                              \label{j:local:mom:oldgap}\\
 -75D/104&\le J_{1,o}\le3D/2.                           \label{j:local:mom:outer}
\end{align}
The first two lines follow from the gap integral identities above,
using $\Gamma\le5/2$ and $J\le\rho+1$. For the outer $J_1$ bounds,
$J_{1,o}\ge-3T_-$, $J_{1,o}\le3J_o\le3(T_++T_-)$ and
\eqref{j:core:gap-defect} give the stated constants.

For the previously defined $P_1,P_2$, $J_{j,c}=mP_j(A,B)$. Both polynomials increase in $B$; their $A$ derivative changes from positive to negative
at $A=B$. A rectangle minimum is therefore at $B_l$ and an $A$ endpoint.
Multiply a negative minimum by an upper bound for $m$, and a nonnegative
minimum by a lower bound. For the upper central charge, maximizing at
$(\bar a,\bar b)$ if $p\le1/2$ gives
$J_{1,c}\le v^2(2-3p)/(6pq)$, a decreasing function of $p$ on that interval.
If $p\ge1/2$, maximize first at $A=B$ to get
$J_{1,c}\le(2/3)v^2q^2/p\le1/3$. Thus every cell may use
\begin{equation}\label{j:local:mom:central-upper}
 J_{1,c}\le\max\{1/3,(2-3l)/(6l(1-l))\}.
\end{equation}

For $|X|\le R$, $R\ge1$, define
\[
 k_4(R)=\frac{(R+1)^2}{R+1/2},\qquad
 k_5(R)=\frac{R^3+2R^2+3R+3/2}{R+1/2}.
\]
Then
\begin{equation}\label{j:local:mom:universal-moments}
 \mu_4\le\min\{R\rho,1+k_4(R)(\rho-1)\},\quad
 \eta_5\le\min\{R\mu_4,1+k_5(R)(\rho-1)\},\quad
 \mu_6\le R\eta_5.
\end{equation}
The first two bounds are the supported moment estimates proved in
Section~\ref{j:fourier:compensated}; the last is pointwise.
The following coarse cancellation estimates retain the radius-six majorants.
The last identity in~\eqref{j:local:re:pricing} gives
\begin{equation}\label{j:local:mom:K-cancellation}
 |K_1+3\rho\beta|\le(96-18\rho)(\rho-1)/13.
\end{equation}
Indeed $u(u+2-3r/2)/(u+1/2)$ is convex on $[0,6]$, its maximum is
$(96-18r)/13$, and its negative part is at most $31/40$, smaller than that
maximum for $1\le r\le37/20$.
The following bounds are used only for $|X|\le6$ and
$1\le\rho\le37/20<2$, with the signed moment hypotheses below.
The other bounds needed on this domain are
\begin{align}
 |K_1|&\le6\rho-3\rho^2/2+3|\beta|,\notag\\
 K_1&\le M+3|\beta|-3\rho g_-/2,&
 -K_1&\le3\rho+3|\beta|+(3\rho/2-3)g_-,\notag\\
 -K_q&\le(\rho/2-3/8)M-\rho/2+1/8+3|\beta|g/2,&
 -K_q&\le8\rho^3/27+(3|\beta|/2-1/2)\rho,               \label{j:local:mom:K-bounds}
\end{align}
where $\mu_4\le M$, $|\gamma|\le g$, $\gamma\ge g_-$, and
$K_q=\eta_5/4-\rho\mu_4/2-3\beta\gamma/2+\rho/2$.
For the first two lines use
$|\E[X|X|^3]-3\gamma|\le3\rho$, and note the negative coefficients
of $\gamma$. For the last line use
$\eta_5\ge\mu_4^{3/2}\ge3\mu_4/2-1/2$, or maximize
$\rho y^2/2-y^3/4$ over $y\ge0$.


\paragraph{Compatible weaker resource constants.}
The implementation also permits the weaker binary charge
$b_I=(616/25)(1-\Theta)^2\operatorname{dist}(p_*,I)^2$.
It follows from $2(6+2\sigma)>616/25$ and
$(1-\theta)^{3/2}\ge(1-\theta)^2$.
With $K=\min(K_+,K_-)$ and $v_* =\max(v_+,v_-)$, the merged charge
\eqref{j:local:re:residual} gives
\[
 T_++T_-\le T_*=(d-b_I)\max\{v_*/[3(A_l+B_l)],(R+2v_*)/K\},
 \qquad d_o=\min\{d,(\sigma+1)T_*\}.
\]
Thus $W_o\le d_o$ and, for $R\le6$,
\[
 -75d_o/104\le J_{1,o}\le3d_o/2,\qquad
 3(\rho J_{1,o}-J_{2,o})\le\frac38
 \max\{\rho^2/(\sigma-1),(R-\rho)^2/(\sigma+1)\}d_o.
\]
These are precisely the merged-cost versions of the sign-specific
resource rule; they do not introduce another available budget.

\paragraph{Conditional lower moments improve two signed charges.}

Let $s_u=\max_{p\in[l,u]}pq$ and $\theta\le\Theta$.
The unsigned cases of~\eqref{j:local:re:conditional} give
\begin{align}
 \mu_4&\ge M_l:=\max\{1,(1-\Theta)^2(1/s_u-3)\},\label{j:local:entry201:m4lower}\\
 \eta_5&\ge E_l:=
 (1-\Theta)^{5/2}\frac{1-4s_u+2s_u^2}{s_u^{3/2}}.\label{j:local:entry201:e5lower}
\end{align}
Both binary moment functions decrease in $s\in(0,1/4]$.
For the fifth this follows because the derivative has the sign of
$s^2+2s-3/2<0$; the fourth is immediate.
The fourth-moment lower bound sharpens the lower endpoint of
$J_1+(\mu_4-3)/6$; its explicit use appears in the cell map below.

The constant signed charge is
\[
 -K_q=\frac{\rho\mu_4}{2}-\frac{\eta_5}{4}
       +\frac32\beta\gamma-\frac{\rho}{2}.
\]
The old coherent estimate used $\eta_5\ge3\mu_4/2-1/2$ and
$\mu_4\le M$. A second estimate uses \eqref{j:local:entry201:e5lower}.
Their difference is independent of $\rho$. Consequently the improved
bound is exactly the old one minus
\begin{equation}\label{j:local:entry201:eta-improvement}
 \frac14\max\{0,E_l-3M/2+1/2\}.
\end{equation}
This is a constant shift on the cell and does not spoil convexity.
It does not assert the generally invalid implication
$\eta_5\ge3M/2-1/2$ from an upper bound on $\mu_4$.

\paragraph{Signed conditional Taylor cancellation.}

Let $K_1=\E[X|X|^3]-3\beta-\tfrac32\rho\gamma$.
On a sign class put $x=|X|$, $a=|V|$, $W=X-V$ and
$\varepsilon=1$ on the positive class, $-1$ on the negative class.
In the actual central-gap domain there is no atom at zero.
With $w_3=x+2a$, $w_4=x^2+2ax+3a^2$ and
$\gamma_V=\E V^3$, exact conditional expansion gives
\begin{equation}\label{j:local:entry201:Kexact}
 K_1(X)-K_1(V)
 =\E\!\left[W^2
 \left\{\varepsilon(w_4-3-\tfrac32\rho w_3)
                   -\tfrac32\gamma_Vw_3\right\}\right].
\end{equation}
Apply~\eqref{j:local:re:conditional} to $j=2,3,4$, with $h=1$
or $h=\varepsilon$, and use
$\rho\gamma-\rho_V\gamma_V
=\rho(\gamma-\gamma_V)+\gamma_V(\rho-\rho_V)$.

The positive coefficient in \eqref{j:local:entry201:Kexact} is
$F(x,a;t)=x^2+2ax+3a^2-3-\tfrac32t(x+2a)$
with $t=\rho+\gamma_V$; the negative coefficient is $-F$
with $t=\rho-\gamma_V$. Their ranges are bounded on complete
rectangles of actual conditional means and support magnitudes.
The positive-support magnitude is at most $R$.
The same formula permits separate actual bounds on the two signed support magnitudes.
On the late domain $p<1/2$, bounds for $\gamma_V$ are
\[
 \sqrt{(1-\Theta)^3(1-2u)^2/[u(1-u)]}
 \le\gamma_V\le(1-2l)/\sqrt{l(1-l)}.
\]
The positive mean lies in
$[\sqrt{(1-\Theta)(1-u)/u},\sqrt{(1-l)/l}]$, and the negative
magnitude mean lies in
$[\sqrt{(1-\Theta)l/(1-l)},\sqrt{u/(1-u)}]$.
Use the lower bound $\rho\ge\rho_B(u)$ only after checking
$3\min(A_l,B_l)\ge(3/2)\sup_I\rho_B$. Indeed conditional expansion
and $(1-\theta)^{3/2}\ge1-3\theta/2$ then imply
$\rho\ge(1-\theta)^{3/2}\rho_B(p)+3\min(A,B)\theta\ge\rho_B(p)$.
Otherwise retain the universal lower bound one.

Since $x+2a>0$, $F$ decreases with $t$.
For fixed $t$ its Hessian is
$\left(\begin{smallmatrix}2&2\\2&6\end{smallmatrix}\right)>0$.
Its rectangle maximum is at a corner. Its minimum is among the four
corners, the minimizer on each edge, and the stationary point
$(x,a)=(3t/8,3t/8)$ if that lies in the rectangle.
On fixed-$x$ edges use $a=t/2-x/3$, clipped to the edge; on
fixed-$a$ edges use $x=3t/4-a$, also clipped.
This is a complete finite extremum list.
Multiplication by $0\le\theta\le\Theta$ uses
$\Theta\min(0,L)$ and $\Theta\max(0,U)$ for the lower and upper
coefficient bounds. It intersects this new interval with the retained
absolute-error interval; neither need be discarded.

\paragraph{The complete central rectangle.}
The central part of $\rho J_1-J_2$ is exactly
$mG(A,B;\rho)$, where $G=\rho P_1-P_2$.
Let a proved probability cell give
$A\in[A_l,A_u]$, $B\in[B_l,B_u]$, $m\in[m_l,m_u]$.
Here
\[
 A_u=\sqrt{u/(1-u)},\quad B_u=\sqrt{(1-l)/l},\quad
 m_l=\sqrt{(1-\Theta)s_0},\quad
 m_u=\sqrt{\max_{p\in[l,u]}pq}.
\]

\begin{lemma}[Complete central rectangle]\label{j:local:entry201:central}
Write $G(A_l+(A_u-A_l)x,B_l+(B_u-B_l)y;\rho)
=\sum t_{kh}x^ky^h$. Define the 25 rational coefficients
\[
 b_{ij}(\rho)=\sum_{\substack{k\le i\\h\le j}}
 t_{kh}\frac{\binom{i}{k}}{\binom4k}
          \frac{\binom{j}{h}}{\binom4h},\qquad 0\le i,j\le4,
\]
and
\[
 \mathcal G(\rho)=\max_{i,j}
       \{m_l b_{ij}(\rho),\,m_u b_{ij}(\rho)\}.
\]
Then $\rho J_{1,c}-J_{2,c}\le\mathcal G(\rho)$.
Moreover $\mathcal G$ is convex in $\rho$.
\end{lemma}
\begin{proof}
Apply Lemma~\ref{j:cert:bernstein} separately to the
two coordinates after the affine substitution. The Bernstein basis is nonnegative and sums to one on the closed square,
so $G$ is a convex combination of its 25 coefficients.
For each coefficient retain both $m$ endpoints, or equivalently use
$m_u$ when it is nonnegative and $m_l$ otherwise. This gives a uniform
bound on the entire three-dimensional rectangle, including boundaries.
Each endpoint-multiplied coefficient is affine in $\rho$; their maximum
is convex. No interior stationary point is omitted.
\end{proof}



\paragraph{One interval-feedback theorem.}
\label{j:local:entry:feedback}\label{j:local:fourier:uniform}

The signed identity from the preceding section is
\begin{equation}\label{j:local:entry:remainder}
 \cR=C\{3\beta J/2+3\rho J_1-3J_2-K_q\}
 +H(CK_1+\varphi U)-H^2(C\rho+\gamma\varphi/2)+E_n,
\end{equation}
where $U=J_1+(\mu_4-3)/6$, $\varphi=\phi(H/\sqrt n)$.
For any one coherent low-frequency triple and complete high-frequency value $K_H$ at floor $N$,
\begin{equation}\label{j:local:entry:fourier-budget}
 |E_n|\le K_F+s_F|H|,\quad
 K_F=c_\pi(L_0+K_H+L_2/\sqrt N),\quad s_F=c_\pi L_H,\quad c_\pi\ge2/\pi.
\end{equation}
The implementation uses $16/25$ before the signed phase and $212/333$
afterward; both are valid since $\pi>333/106$.

\begin{lemma}[Continuous cells and feedback]\label{j:local:entry:cell-theorem}\label{j:local:entry:corners}
Suppose bounds for one actual cell give
\[
 r_0\le\rho\le r_1,\quad0<c_0\le C\le c_1,\quad
 |\beta|\le b,\quad |\gamma|\le g,\quad \mu_4\le M,\quad D\le d,
\]
$K_1\in[k_l,k_u]$, $U\in[u_l,u_u]$, and
$0\le\varphi_-\le\varphi\le2/5$. Freeze the shared outer charge $d_o$,
lower fifth moment $E_l$, upper $s_u=\max_{p\in[l,u]}p(1-p)$, and
central envelope $\mathcal G$ from the preceding bounds.
Define
\begin{align*}
 A(r)&=\tfrac32 b(r+2\sqrt{s_u})+3\mathcal G(r)
 +\tfrac38\max\{r^2/(\sigma-1),(R-r)^2/(\sigma+1)\}d_o\\
 &\quad +(r/2-3/8)M-r/2+1/8+\tfrac32bg
 -\tfrac14\max\{0,E_l-3M/2+1/2\},\\
 B(c)&=\max_{k\in\{k_l,k_u\},\,u\in\{u_l,u_u\},\,
                      \psi\in\{\varphi_-,2/5\}}|ck+\psi u|,\\
 Q(c,r)&=cr+\tfrac12\varphi_-(3.16r-3-d)
       -\tfrac12(2/5-\varphi_-)(d-.16)_+.
\end{align*}
If all four corner values of $Q$ are positive, then, for every real $H$,
\begin{align}
 \cR&\le\max_{(c,r)\in\{c_0,c_1\}\times\{r_0,r_1\}}
       \{K_F+cA(r)+(B(c)+s_F)|H|-Q(c,r)H^2\}\notag\\
 &\le\max_{c,r\text{ corners}}
       \{K_F+cA(r)+(B(c)+s_F)^2/(4Q(c,r))\}.           \label{j:local:entry:four-corners}
\end{align}
For the following feedback, assume in addition that the cell contains a
ratio extremizer with $e=C_n-\CE\ge0$, $n\ge N$, and that
$0<T<.4N$ is a strict upper bound for the final maximum in
\eqref{j:local:entry:four-corners}. Then $T<nC_n\rho$, since
$C_n\rho\ge\CE>.4$, and the weight-removal step in
Corollary~\ref{j:core:license} gives
\begin{equation}\label{j:local:entry:feedback-values}
 D<\frac{15000T}{997N},\quad
 \varphi>\frac{997}{2500}\left(1-\frac{5T}{2N}\right),\quad
 C<\CE+V/N\quad\text{if also }\cR/\rho<V.
\end{equation}
Only positive density floors are used. These statements remain valid at
$H=0$, on closed cell edges, and for zero conditional variance.
\end{lemma}
\begin{proof}
The constant coefficient follows from the shared outer budget,
Lemma~\ref{j:local:entry201:central}, and~\eqref{j:local:mom:K-bounds}, using
\[
 J\le\rho+2ABm\le\rho+2(1-\theta)^{3/2}\sqrt{pq}
       \le\rho+2\sqrt{s_u}.
\]
This follows from $A\le\bar a$, $B\le\bar b$, $m^2=(1-\theta)pq$;
it is used only in the charge $3|\beta|J/2$. Since $\gamma\ge3.16\rho-3-d$, its negative
part is at most $(d-.16)_+$. Using the upper density only for that negative
part gives $C\rho+\gamma\varphi/2\ge Q$. Absolute value is separately
convex in each coordinate of $ck+\psi u$, proving the formula for $B$.
For fixed $r$ and $|H|$, the resulting upper bound is convex in $c$.
Take those endpoints \emph{before} completing the square. At either
endpoint $c>0$, it is convex in $r$, since $\mathcal G$ and the maximum of the two outer quadratics are
convex, every other term of $A$ is affine, and $Q$ is affine in $r$. This proves the first line of~\eqref{j:local:entry:four-corners};
$fh-qh^2\le f^2/(4q)$ proves the second. Bilinearity of $Q$ proves
its positivity throughout the rectangle from its corners.

\end{proof}

\paragraph{Signed phase curvature in the resource schedule.}
\label{j:local:re:signed-phase-repair}
The auxiliary Fourier orientation does not change the signed third moment
in the actual transfer identity. The paired phase estimate must therefore
retain its sign. On a fixed resource cell put
\[
 L=\frac{997}{2500}\left(1-\frac{5r}{2N}\right),\qquad
 g_0(\rho)=3.16\rho-3-d,\qquad U=.4,
\]
where $r,d$ are the preceding raw and defect upper bounds, and the
checker requires $0\le L\le U$. The actual density lies in $[L,U]$
and the actual signed moment satisfies $\gamma\ge g_0(\rho)$. Consequently
\begin{equation}\label{j:local:re:sign-aware-price}
 C\rho+\tfrac12\gamma\phi
 \ge q_{\rm s}(C,\rho):=
 C\rho+\tfrac12\min\{L g_0(\rho),U g_0(\rho)\}.
\end{equation}
This follows by first multiplying the lower bound on $\gamma$ by the
nonnegative actual density, and then minimizing over its two endpoints.
The simpler expression $C\rho+Lg_0(\rho)/2$ is licensed if $g_0\ge0$,
or if nonnegativity of the actual $\gamma$ has been separately proved.
It is not licensed merely by reflecting the auxiliary Fourier integral.

The corrected price preserves the four-corner reduction. For fixed $C$,
$q_{\rm s}$ is concave in $\rho$, being affine plus the minimum of two
affine functions. For fixed $\rho$ it is affine in $C$. In the phase
envelope $A+f|H|-q_{\rm s}H^2$, the remaining terms are separately
convex in $C,\rho$, by the moment and coefficient bounds already proved.
Thus the entire envelope is separately convex, and two successive
one-dimensional endpoint bounds place its maximum at the four corners.
The same argument applies after adding the affine conserved-resource
curvature $e_\rho$ and subtracting a barrier price $K\rho$.
If the four corrected divisors are positive, concavity makes them
positive throughout the rectangle, and square completion gives the
valid raw and resource bounds
\[
 \max_{C,\rho\text{ corners}}\left(A+\frac{f^2}{4q_{\rm s}}\right),
 \qquad
 \max_{C,\rho\text{ corners}}
       \left(A+\frac{f^2}{4(q_{\rm s}+e_\rho)}-K\rho\right).
\]
Any threshold-root export likewise uses the corrected corner curvature.

For the later acquired domains, the conditional cubic bound
\eqref{j:local:re:gamma} applies in the actual contact orientation.
It requires the same marginal's probability interval, support radius
and variance ceiling, all established before the later paired calls.
The exact domain check covers 302 input rows: 45 central-input rows,
45 left-input rows, 45 first-confinement feedback rows, 45 Sperner feedback
rows, 84 signed-resource rows, and 38 final left-radius rows. With the
common already acquired radius $R=3.84$, each right side is greater than
$.096$. Thus all later paired calls have an actual nonnegative signed
third moment, independently of any auxiliary Fourier orientation.

There is also a convenient cell-constant correction. Since $\rho\ge1$,
$(-g_0(\rho))_+\le(d-.16)_+$, so the weaker price
\[
 C\rho+\tfrac12L g_0(\rho)
       -\tfrac12(U-L)(d-.16)_+
\]
is valid. This is the correction selected for the additional index--46
calculation. Both forms keep the raw remainder separate from the
barrier-subtracted resource budget.

The broad preliminary and complementary resource passes are replayed with
\eqref{j:local:re:sign-aware-price}; their subsequent consumers retain the same
fixed target tables only when the complete corrected replay verifies every
exported resource. A diagnostic equality of two corner maxima alone does
not replace this check, since threshold bounds also consume the phase
coefficients. Later uses of the simpler formula require their own explicit
nonnegative signed-moment or small-defect license.




\paragraph{Using a floor calculation at a later index.}
\label{j:local:transfer}
A floor calculation encloses one fixed attaining ratio marginal with
$C\ge\CE$; it does not assert monotonicity of $C_n$. For actual
$n\ge N$, every saved pass
preserves the two simultaneous statements
\begin{equation}\label{j:local:ext:floor-invariant}
 \mathcal R_n\le r,\qquad nD\le Nd.
\end{equation}
In particular $D\le d$ may be used in index-free moment and gap
inequalities. The stronger weighted statement is retained whenever a
cost is proportional to the sample size.

Here is why each operation in the selected schedules preserves this
invariant. Keep the actual marginal's moments, support, signed phase
and admissible frequency windows fixed while increasing $n$.
All numerical gates mentioned below are checked before the operation.

\emph{Fourier bounds.} Positive sums and products of nonnegative
decreasing factors remain decreasing. For the complete HIGH bound
of Proposition~\ref{j:fourier:high}, these factors are
\[
 \sqrt n e^{-cn}\ (2Nc>1),\quad n^{-1},\quad
 e^{-knb^2}/\sqrt n,\quad
 r(x)=\frac{x+2}{x^2+4x+2},\quad x\text{ increasing in }n.
\]
Indeed $r'=-(x^2+4x+6)/(x^2+4x+2)^2<0$.
The comparator $\sqrt n e^{-n/8}\{2r(n/8)+g_c/3\}$ decreases
for $N\ge5$. The improved first piece has the same factors with
its separately checked positive rate and slope. For the optional zero
piece, the derivative of $\sqrt{x}e^{-x}r(x)$ has the sign of
\[
 4-10x-22x^2-13x^3-2x^4.
\]
This polynomial decreases for $x>0$ and equals $-11/128$ at
$x=1/4$, so the native gate $Nkb^2\ge1/4$ suffices. Every use
retains the event/window conditions $0<a\le1$, $m_e>0$,
$0<d_0\le12$, $1\le M\le111/10$, $0<w\le1/4$ and positive
sine slope. Taking the minimum of complete admissible windows is
valid even when the minimizing window changes with $n$.

For LOW, the nonconstant factors are
\[
 (n/(n-1))^{3/2},\quad (n/(n-1))^{3/2}/(n-1),\quad
 a_n^{-j},\quad b_n^{-j},\quad n^{-1/2},
\]
where $a_n=(n-1)/(2n)$, $b_n=(45-M)(1-2/n)/96$ and $j>0$.
They decrease on the stated positive domains. This covers the full
error $e/(2a_n^2)+B^2/(2b_n^3)$ and $L_2/\sqrt n$; the paired
moment-rectangle hypotheses still have to hold.

\emph{Density and score updates.} Under the raw sign and positivity
license $r<.4N$, the invariant \eqref{j:local:ext:floor-invariant} gives
\[
 \phi(H/\sqrt n)\ge\phi_{0,l}(1-5r/(2n)).
\]
Its lower bound increases with $n$, whereas
$r/n-\phi_{0,l}(1-5r/(2n))D_l/6$ decreases for $r,D_l\ge0$.
Both may therefore be evaluated at $N$. The general curvature
correction has derivative $(g_0+t)/2\ge0$ with respect to this
density floor, where $g_0=(79/25)\rho-3-d$, $t=(d-4/25)_+$:
if $d\le4/25$, use $g_0\ge4/25-d$; otherwise
$g_0+t=(79/25)(\rho-1)$. The signed curvature
\[
 C\rho+\tfrac12\min\{\phi_lg_0,(2/5)g_0\}
\]
also increases with $\phi_l$ when $g_0\ge0$ and is constant
when $g_0<0$. Thus both signs are retained. Positive added
curvatures and all divisors are checked before square completion.

For $C\ge\CE$ and $\CE<40974/10^5$, the compulsory charge is
\[
 n(C-\CE)\rho\ge
 N\max\{0,C-40974/10^5\}\rho_l.
\]
Split the score cover at $40974/10^5$; the charge is affine on
each side. The coherent corner calculation then exports
$nD\le B/\kappa$, hence $d=B/(\kappa N)$, for
$\kappa=997/15000$. A negative budget excludes the cell before
division. Positive $K/n$ calculations separately retain their native
$K\rho$ charge; two lower bounds on the same excess are not added.
Closed covers, upward rounding and intersection with an earlier
bound now preserve \eqref{j:local:ext:floor-invariant} by finite
induction, without a fixed-point or limiting argument.

\emph{Weighted physical resources.} If an acquired inequality is
\[
 D\ge b+g d_{\rm in}+(N/n)h_N\lambda_n,
 \qquad \lambda_n=ns,\quad b,g,h_N\ge0,
\]
then the invariant gives
$(n/N)(b+g d_{\rm in})+h_N\lambda_n\le d$, hence
$b+g d_{\rm in}+h_N\lambda_n\le d$.
This retains the original joint intercepts. The actual $n$ is still
used in every binomial probability. The transfers used here are
$N=25$ to $26\le n\le30$ and $N=31$ to $32\le n\le34$;
each actual index has its own verified support calculation.
Index-free support refinements use only already acquired data. For
example $6mAB\ge P>0$ gives $(A+B)^2\ge2P/(3m_u)$ by
$4AB\le(A+B)^2$; it does not establish an initial central gap.

\paragraph{Phase and variance exports from the same budget.}
Suppose one signed phase family gives, for $h=|H|\ge0$,
\[
 0\le D\le c_n\max_j(A_j+F_jh-Q_jh^2),\qquad c_n>0,
 \quad Q_j>0.
\]
Some corner is nonnegative at the actual $h$. For that corner
$\Delta_j=F_j^2+4Q_jA_j\ge0$ and
$h\le(F_j+\sqrt{\Delta_j})/(2Q_j)$.
If $u_j\ge0$ is a rational upper enclosure of $\sqrt{\Delta_j}$,
the complete family therefore gives
\[
 h\le\max_j r_j,\qquad
 r_j=\begin{cases}0,&\Delta_j<0,\\
 \max\{0,(F_j+u_j)/(2Q_j)\},&\Delta_j\ge0.
 \end{cases}
\]
This includes double roots and $h=0$. Keep separate maxima for
the two phase signs and intersect each with its previous cap.
At floor31 the normalization is $15000/(997\cdot31)$; its
actual-index version $c_n=15000/(997n)>0$ leaves this argument
unchanged. Positivity of every $Q_j$ is required.

Finally suppose bounded conditional supports imply $s\le B\delta$.
Writing $O_0,O_1$ for their outward offsets, take
$B=\max(O_0,O_1)\max(p_u,1-p_l)>0$.
Combining this with the original joint
plane $\delta/d_0+\lambda_n/L\le1$ and $\lambda_n=ns$ gives
\begin{gather*}
 \lambda_n\le\{L^{-1}+(nBd_0)^{-1}\}^{-1},\\
 \theta_n=\frac{\lambda_n}{npq+\lambda_n}
 \le\frac1{1+n v_l/L+v_l/(Bd_0)}
 \le\frac1{1+N v_l/L+v_l/(Bd_0)},
\end{gather*}
where $v_l=\min_{p\in[p_l,p_u]}pq>0$ and $d_0,L>0$.
Upward rounding and intersection with the previous variance bound
preserve this conclusion. A downward support-scale enclosure enlarges
$O_0,O_1$, so also weakens the bound safely. The refinement concerns
physical $\theta_n$; it does not identify $\lambda_n$ with a
floor's diagnostic $\lambda_N$ or replace the plane's intercept $L$.
All 193 saved first-band refinements at floor 25 are checked at each
actual index 26--30.


Here is the version of the moment matrices that permits different
positive and negative support radii. On either sign write $Y=|X|$,
$0\le Y\le r$, and let $p_s,m,a,b,c_4,c_5$ be its unnormalized
moments of orders zero through five. Positivity of the weights
$r-Y$, $Y(r-Y)$, and $Y^2(r-Y)$ gives
\[
 \begin{pmatrix}rp_s-m&rm-a\\rm-a&ra-b\end{pmatrix},\quad
 \begin{pmatrix}rm-a&ra-b\\ra-b&rb-c_4\end{pmatrix},\quad
 \begin{pmatrix}ra-b&rb-c_4\\rb-c_4&rc_4-c_5\end{pmatrix}
 \succeq0.
\]
For positive denominators their determinant inequalities imply
\begin{align*}
 p_s&\ge m/r+(rm-a)^2/[r(ra-b)],\\
 c_4&\le rb-(ra-b)^2/(rm-a),\\
 c_5&\le rc_4-(rb-c_4)^2/(ra-b).
\end{align*}
A zero diagonal entry forces the corresponding off-diagonal entry to
vanish. One may check that condition or discard the nonnegative
correction, but one cannot divide by zero. In interval form a correction
is bounded below by the square of a positive-part lower numerator
divided by an upper positive denominator. The sign radius $r$ is used
only for its own class. Summing the first inequality over the two signs
and using $p_++p_-\le1$ can improve the common upper mean.

One further bound used at small indices strengthens the signed constant
coefficient. Set $\delta=\gamma-\rho\beta$ and $Y=|X|$. Then
\[
 V:=\E\{Y(Y^2-\rho Y)^2\}
   =\E|X|^5-2\rho\mu_4+\rho^3,
 \qquad
 \delta=\E\{Y\operatorname{sgn}(X)(Y^2-\rho Y)\}.
\]
Here $m>0$ by centering and unit variance.
Cauchy's inequality for the nonnegative measure $Y\,dP$ gives
$\delta^2\le2mV$. Consequently
\begin{equation}\label{j:local:ext:weighted-cubic-loss}
 \frac\rho2\mu_4-\frac14\E|X|^5
 \le\frac{\rho^3}{4}-\frac{(\gamma-\rho\beta)^2}{8m}.
\end{equation}
This deduction can be combined with the two other deductions in
Section~\ref{j:local:atlas} by taking their maximum. Their sum would spend
the same nonnegative quantity more than once.



\paragraph{A universal lower bound for each sign probability.}
\label{j:local:ext:sign-mass}
Six probability panels at index six are excluded before the count
calculation. The exclusion is the following elementary moment inequality;
it does not require support, a central gap, or absence of a zero atom.

\begin{lemma}[Sign mass and third moment]\label{j:local:sign-mass}
If $\E X=0$, $\E X^2=1$, and $\rho=\E|X|^3<\infty$, then
\[
 \rho^2\Pp(X>0),\quad \rho^2\Pp(X<0)
 \ \ge\ \frac{8(7\sqrt7-17)}{27}
 \ >\ \frac{11261}{25000}.
\]
\end{lemma}
\begin{proof}
Write $p=\Pp(X>0)$, $a=\E X_+^2$, and
$m=\E X_+=\E X_->0$. Centering and unit variance imply
$0<a<1$. Cauchy's inequality, first with factors
$X_+^{1/2}$ and $X_+^{3/2}$ and then with their negative-sign
counterparts, gives
\[
 a^2\le m\E X_+^3,\qquad (1-a)^2\le m\E X_-^3.
\]
Also $m^2\le pa$. Consequently
\[
 p\rho^2\ge f(a):=\frac{(a^2+(1-a)^2)^2}{a}.
\]
The derivative has the sign of $6a^2-2a-1$, because
\[
 f'(a)=\frac{(2a^2-2a+1)(6a^2-2a-1)}{a^2}.
\]
Thus the unique minimum on $(0,1)$ is at
$a=(1+\sqrt7)/6$, where
$f(a)=8(7\sqrt7-17)/27$. The strict rational comparison follows
from $2645751^2<7\cdot10^{12}$ and
$8(7\cdot2645751/10^6-17)/27>11261/25000$.
Apply the same argument to $-X$ for the other sign probability.
\end{proof}



\paragraph{Scalar consequences of the common inward budget.}
\label{j:local:ext:inward-scalar}
Some finite count domains use a stronger absolute-third-moment floor or
noise ceiling before their five-dimensional subdivision. The following
calculation supplies those scalar restrictions. It uses the sign
conditioning of \eqref{j:core:conditional-moments}.

Reflect only for the scalar calculation, so that
$0<p\le1/2$, $q=1-p$, $v=pq$. This preserves $\rho$, total
conditional variance $s$ and total inward displacement. The count
probability, its phase and its original orientation are not reflected.
Suppose the inward endpoints obey $a_0+a_1\le d$ with $0<d\le p$.
The conditional third-moment inequalities and Cauchy's inequality give
\begin{equation}\label{j:local:ext:inward-I}
 \rho\ge I(p,s,d):=
 \frac{v(1-2v)+(3p-d)s+s^2/(qd)}{(v+s)^{3/2}}.
\end{equation}
Indeed, writing $b=qs_0$, $c=ps_1$ and $b+c=s$, the absolute-cubic
numerator is at least
\[
 v(1-2v)+3pb+3qc-a_0b-a_1c
                       +\frac{b^2}{qa_0}+\frac{c^2}{pa_1}.
\]
Its linear part is at least $(3p-d)s$, while its quadratic part is
at least $(b+c)^2/(q(a_0+a_1))\ge s^2/(qd)$.
Vanishing inward endpoints force the corresponding centered variance
to vanish; those terms are handled before division. If $d=0$, then
$s=0$ directly.

The derivative has the exact numerator
\begin{equation}\label{j:local:ext:inward-I-derivative}
 (v+s)^{5/2}I_s
 =v(6p-3p^2-d-3/2)
   +\left\{\frac{2p}{d}-\frac{3p-d}{2}\right\}s
   +\frac{s^2}{2qd}.
\end{equation}
Consequently, on $p\in[l,u]\subset(0,1/2]$, the two gates
$d\le l$ and $f:=6l-3l^2-d-3/2\ge0$ imply
\[
 \rho\ge\rho_B(p):=\frac{1-2p(1-p)}{\sqrt{p(1-p)}}
                  \ge\rho_B(u).
\]
The last function is decreasing on this half-interval. If a previously
acquired ceiling is $\rho\le R$, a point with
\[
 R^2p(1-p)-(1-2p(1-p))^2<0
\]
is impossible. This polynomial increases on $0<p\le1/2$, because
its derivative is $(1-2p)(R^2+4-8p(1-p))\ge0$.
Exact bisection may therefore raise the closed lower probability
endpoint to a known infeasible point; retaining that endpoint is a
conservative enlargement.

When $f>0$ and $s\le S$, formula~\eqref{j:local:ext:inward-I-derivative}
also gives
\[
 I_s\ge
 \frac{l(1-l)f}{[u(1-u)+S]^{5/2}}.
\]
Integration from zero supplies a new variance ceiling. Directed square
roots and outward division make every finite update valid separately.

A useful weaker floor survives failure of the monotonicity gate.
Assume still $0<d\le l$ and set
\[
 a=(3/2-6l+3l^2+d)_+,
 \qquad B=2l/d-3u/2\ge5/4.
\]
The only possible decrease of $I$ is bounded below by
$-(va-Bs)_+/v^{5/2}$. Integrating this negative triangle gives
\begin{equation}\label{j:local:ext:inward-deficit}
 \rho\ge\rho_B(u)-\frac{a^2}{2B\sqrt{l(1-l)}}.
\end{equation}
In particular, failure of the monotonicity gate must not be replaced
by the unsupported assertion $\rho\ge\rho_B(p)$.

There is a second variance contraction which works with either sign
of $f$. The global Taylor upper bound for $(v+s)^{3/2}$, combined
with \eqref{j:local:ext:inward-I}, gives the necessary inequality
\begin{equation}\label{j:local:ext:scalar-noise-quadratic}
 Gs^2+fs\le\epsilon\,[u(1-u)+S]^{3/2},
\end{equation}
where
\[
 \epsilon=(R-\rho_B(u))_+,\qquad
 G=\frac1{(1-l)d}-\frac{3(1-2l(1-l))}{8l(1-l)}
       \ge\frac5{8l(1-l)}+\frac34>0.
\]
To see this, subtract $\rho_B(p)(v+s)^{3/2}$ from the lower
numerator in \eqref{j:local:ext:inward-I}, use the quadratic Taylor upper
bound on the subtracted term, and then use $p\ge l$ and
$\rho_B(p)\ge\rho_B(u)$. The positive part in $\epsilon$ is
necessary when $f<0$. The larger real root of
\eqref{j:local:ext:scalar-noise-quadratic}, enclosed from above, is a valid
new upper bound for $s$, even for negative $f$.

Finally suppose the original joint resource has the stronger form
\[
 d_{\rm actual}\le d_0(1-s/S_0),\qquad S_0=\Lambda/n.
\]
To test a proposed upper endpoint $a\in(0,S_0)$, assume $s\ge a$.
Then $d_{\rm actual}\le d_0(1-a/S_0)$, so the preceding quadratic
calculation can be repeated with this smaller inward bound. If its
resulting upper bound is strictly less than $a$, the whole region
$s\ge a$ is impossible. Otherwise the probe proves no lower bound
on $s$. A finite sequence of such hypothetical tests is sound
without any convergence assertion. Throughout it, the original
$S_0$ remains in the joint budget; a separately improved noise
ceiling clips that budget and does not replace its intercept.


\subsection{Paired Fourier cancellation}
\label{j:local:paired}
The triangle LOW bound separates the fourth- and fifth-order errors.
Near a two-point law their leading imaginary parts partly cancel.
We retain that cancellation in one Gaussian norm; all other Taylor and
power remainders remain those of Proposition~\ref{j:fourier:low}.

\begin{lemma}[The cancellation norm and its finite tail bounds]
\label{j:local:entry189:mills}
Put $I=\sqrt{\pi/2}$ and
\[
 J(z,c)=\int_0^\infty e^{-x^2/2}|zx^2-cx^4|\,dx.
\]
The function $J$ is convex in $(z,c)$ and invariant under simultaneous
reflection. Its maximum on a rectangle is therefore bounded by its
largest corner value. For $zc\le0$ it equals $I(|z|+3|c|)$.
For $z,c>0$, put $x=z/c$ and define
\[
 d_{k+1}=1,\qquad d_j=1+\frac{j}{xd_{j+1}}\quad(j=k,\ldots,1),
 \qquad m_k(x)=d_1^{-1}.
\]
Then $J$ is bounded by each of
\begin{align}
 U(z,c)&=I|z-3c|+6\sqrt{zc}\,e^{-x/2},\notag\\
 V(z,c)&=I(z-3c)+\sqrt{zc}\,e^{-x/2}
       \left(6+\frac{2(3-x)}x m_k(x)\right),
 \label{j:local:entry189:mills-bound}
\end{align}
where $k=11$ for $x\ge3$ and $k=12$ otherwise. The parenthesized factor
is positive. For negative $z,c$, apply these formulas to $|z|,|c|$.
\end{lemma}
\begin{proof}
Convexity follows by integrating the absolute value of a linear function;
successive endpoint bounds give the rectangle assertion. For $zc>0$,
reflection and integration by parts, split at $r=\sqrt{z/c}$, give
\[
 J=I(z-3c)+6cr e^{-r^2/2}
       +2(3c-z)\int_r^\infty e^{-t^2/2}\,dt.
\]
Dropping a nonpositive last term, or bounding its integral by $I$ when
positive, proves $U$. To retain the signed tail, set
$I_j(r)=\int_0^\infty t^j e^{-rt-t^2/2}\,dt$ and
$q_j=I_j/I_{j-1}>0$. Integration by parts yields
\[
 I_0=\frac1{r+q_1},\qquad q_j=\frac j{r+q_{j+1}},\qquad
 I_0=e^{r^2/2}\int_r^\infty e^{-t^2/2}\,dt.
\]
Truncating at $q_{k+1}=0$ reverses the inequality at every reciprocal:
odd truncations are lower bounds, even ones upper bounds. At $r=\sqrt x$
their values multiplied by $r$ are exactly $m_k(x)$ above. The tail's
multiplier $3c-z$ selects the stated parity and proves $V$, without
requiring convergence of an infinite fraction. For $x<3$ the
parenthesized factor exceeds $6$; for $x\ge3$, $0<m_{11}(x)<1$ makes it
greater than $6-2(x-3)/x>4$.
\end{proof}

For outward evaluation use an upper enclosure for $I$ in $U$ and in
the case $zc\le0$. In $V$, use a lower enclosure when $z-3c<0$
and an upper one otherwise; the remaining factor is positive. Thus
square roots and $e^{-x/2}$ may be bounded upward. Let $J_u$ be the
largest such corner bound, taking the minimum of $U,V$ when applicable.

\paragraph{A single transfer formula.}
Retain the hypotheses and notation of Proposition~\ref{j:fourier:low}:
$N>2$, $1\le\mu_4\le M\le11.1$, $|\gamma|\le g\le1.85$, and actual
ceilings $|a|\le a_*$, $|z|\le Z$, $\mu_6\le M_6$, where
\[
 a=(3-\mu_4)/24,\qquad z=(10\gamma-\E X^5)/120,
 \qquad c=a\gamma/6.
\]
Any rectangle containing the actual pair $(z,c)$ may be used above.
For $F_N\ge(N/(N-1))^{3/2}$ and $I_u\ge I$, a valid LOW triple is
\begin{align}\label{j:local:fourier:paired}
 L_0&=F_NJ_u+\frac{F_NI_u}{N-1}
       (35a_*g^3/48+35Zg^2/24)+251g^3/17280,\\
 L_H&=F_NI_u\{a_*+5Zg/[2(N-1)]\}+251g^2/4800,\notag\\
 L_2&=E/(2a_0^2)+B^2/(2b_0^3),\notag
\end{align}
with the same $E=M_6/720+g^2/72+(1/2+g^2/144)^3/5$,
$B=a_*+Z/2+E/4$, $a_0=(N-1)/(2N)$, and any
$0<b_0\le(45-M)(1-2/N)/96$.

Indeed, write $b=n-1$ and $v=b\gamma t^3/6$. The leading imaginary
part in the common Taylor proof, after division by $t^3$ and rotation
by $Ht$, is bounded by
\[
 |zt^2-atv|+a_*t|\sin v-v|+Zt^2|\cos v-1|
       +|H|t\{a_*t+Zt^2|\sin v|\}.
\]
Multiply by $ne^{-bt^2/2}$ and integrate. The first term contributes
$nJ(z,c)/b^{3/2}$; the others use
$|\sin v-v|\le|v|^3/6$, $|\cos v-1|\le v^2/2$ and
$|\sin v|\le|v|$. The recurrence
$\int_0^\infty x^{2j}e^{-x^2/2}\,dx=(2j-1)!!I$
gives respectively the coefficients $35/48,35/24,1,5/2$ in
\eqref{j:local:fourier:paired}. The common power/Taylor errors and
the signed Gaussian comparator give exactly the unchanged $L_2$ and
the two $251$ terms. All inverse-index factors decrease for $n\ge N$;
the conclusion is uniform when the law and threshold vary with $n$.

\paragraph{One moment envelope for the coefficient rectangle.}
The broad triangle choice is already proved in
Section~\ref{j:fourier:pearson-envelope}. With its exact radius-six
and actual-radius quantities $Z_0,Z_r$, it uses
\begin{equation}\label{j:local:fourier:z0}
 120Z=\min\{E_5+10g,Z_0,Z_r\},\qquad
 a_*=\max(2,M-3)/24,\quad M_6=RE_5.
\end{equation}
Here $E_5\ge\E|X|^5$ and $|X|\le R\le6$ are previously acquired.
The radius-six derivative test and the actual-radius clipped vertex
are exactly those in that proof; no moment ceiling is substituted
from an unrelated marginal. A positive rectangle gives the stronger
paired choice.

\begin{proposition}[A positive cancellation rectangle]
\label{j:local:paired-rectangle}
Suppose $\E X=0$, $\E X^2=1$, $|X|\le R$, $1\le\mu_4\le M<3$,
$0\le\gamma_l\le|\gamma|\le g$, $g^2<2$, $R>2g$, and
$M-1-\gamma_l^2\ge0$. After reflecting both odd Fourier coefficients
if necessary, set
\begin{gather*}
 t_* =\min\{g,(10-2M)/(2(R-g))\},\\
 120Z=\min\{E_5+10g,R(M-1)+(10-2M)t_*-(R-g)t_*^2\},\\
 z_l=\max\{-Z,[8\gamma_l-\gamma_l^3
        -(R+2\gamma_l)(M-1-\gamma_l^2)]/120\},\quad z_u=Z,\\
 a_*=(2-\gamma_l^2)/24,\qquad
 c_l=(3-M)\gamma_l/144,\quad c_u=(2-\gamma_l^2)g/144,\\
 M_6=\min\{RE_5,1+3g^2+g^4+
 (R^2+2Rg+3g^2+2)(M-1-\gamma_l^2)\}.
\end{gather*}
These are valid coefficient bounds and $[z_l,z_u]\times[c_l,c_u]$
contains $(z,c)$. If $3g^2\le2$, one may further decrease $c_u$ to
$\min\{c_u,(2g-g^3)/144\}$. Use the damping constant
$b_0=(45-M)(1-2/N)/96$; for $M\le2$ the weaker
$43(1-2/N)/96$ is also valid.
\end{proposition}
\begin{proof}
For $Q=X^2-\gamma X-1$, the common Pearson identities and one further
multiplication give
\begin{align}\label{j:local:fourier:pearson}
 \E Q^2&=\mu_4-1-\gamma^2,&
 \E X^5&=\gamma^3+2\gamma+\E[Q^2(X+2\gamma)],\notag\\
 \mu_6&=1+3\gamma^2+\gamma^4+
       \E[Q^2\{(X+\gamma)^2+2\gamma^2+2\}].
\end{align}
Thus $\mu_4\ge1+\gamma^2$, proving $a_*,c_l,c_u$; the optional
upper bound uses increasing $2t-t^3$ on $[0,g]$.
The positive-kernel estimates of Section~\ref{j:fourier:pearson-envelope}
give $-Q_R(t)\le120z\le P_R(t)$ for $t=\gamma\ge0$.
Since $M<3$, their absolute upper envelope is
$R(M-1)+(10-2M)t-(R-g)t^2$, maximized at $t_*$.
Moreover $(-Q_R)'(t)=10-2M+2Rt+3t^2>0$, proving $z_l$.
In the last identity, bound each nonnegative factor upward using
$t\le g$ and $\E Q^2\le M-1-\gamma_l^2$ to obtain $M_6$.
The transfer's domain follows from $M<3<11.1$, $g<\sqrt2<1.85$.
\end{proof}

Finally a probability row must supply a lower bound for the actual
$\rho$, not automatically its Bernoulli counterpart. If
$p\in[l,u]\subset(0,1/2)$, $\theta\le\Theta<1$, and acquired inner
endpoints are at least $a\ge0$, the conditional cubic identity gives
$\rho\ge(1-\theta)^{3/2}\rho_B(p)+3a\theta$.
For outward $r_-\le\rho_B(u)$, $r_+\ge\rho_B(l)$ this implies
\[
 \rho\ge r_*:=\max\{1,r_--(3r_+/2-3a)_+\Theta\},
\]
because $\rho_B$ decreases on $(0,1/2)$ and
$(1-\theta)^{3/2}\ge1-3\theta/2$.
The same $r_*$ is retained in all Fourier corners, signed phase prices,
conserved budgets and the signed fourth-moment errors
\eqref{j:local:entry201:Kexact}; the binary-$K$ alternative still requires
$1-\Theta>2/3$. Reflection here concerns only the auxiliary norm and
does not change the actual contact orientation.

\subsection{Confinement from the remaining defect budget}
Two different support implications are useful. Gaussian contact excludes
a point by a strict negative upper bound on its contact expression;
closed cells cover the finite exterior and a derivative-sign certificate
covers any remaining ray. The shared Gaussian support theorem supplies
this callback, including the bounded-threshold radius and coordinate-strip
contracts; it is a prior theorem, not a test on moments alone. Noise confinement proves $|U|\le a$ on an
already acquired row. It then implies, without a small-noise limit,
\begin{equation}\label{j:local:re:radiusfeedback}
 |X|\le R_*:=
 \frac{\max\{u+a,1-l+a\}}{\sqrt{\min_Ip(1-p)}}.
\end{equation}
The denominator uses $\E U^2\ge0$, so zero noise is included.
The Fourier/moment formulas remain valid for $R\ge1$; the paired Pearson
choice additionally checks $R>2g$. In particular there is no intrinsic
lower radius $2.3$: the radial ratios in Lemma~\ref{j:local:re:pricing} increase
on all $x\ge0$, and the other formulas use their explicit domains.

Here is the selected proof schedule. Each row names \emph{an application}
of Theorem~\ref{j:local:acquisition}, not a different analytic theorem.
\begin{center}\small
\begin{tabular}{p{.13\linewidth}p{.32\linewidth}p{.45\linewidth}}
\toprule
Node&Input and operation&Output consumed later\\\midrule
I&$N=47$, $R=3.84$, $C\le.4247$; 13 signed passes plus export&
100 closed probability cells; 35 survive; $\mathcal R<1.885$, $nD<47(.590)$\\
C&Central preparation, contact radius $2.5$, probability refinement&
45 cells $[.38,.47]$, width $.002$; canonical simultaneous input\\
L&Complete complementary cover, contact radius $3$, refinement&
39 cells $[.302,.38]$; all other complementary cells excluded\\
C1&First confinement from C; radius formula; enclosure&
First improved central table\\
C2&Antichain confinement from C; intersect C1; enclosure&
Selected central table: 45 cells, $[.38,.47]$\\
S&Separate sign-resource pricing on C2 and L&
One left cell excluded; 38 left cells remain\\
L2&Left contact radius $2.7$ from S; enclosure&
Selected left table: 38 cells, $[.304,.38]$\\\bottomrule
\end{tabular}
\end{center}
Every shared endpoint belongs to its adjacent closed cells. The
confinement producing C2 consumes C, not C2. The selected central local
comparison uses C2; it does not need the optional further central
refinement in S. The final cover is exactly $[.304,.47]$.

\begin{corollary}[Uniform acquired entry]\label{j:local:re:acquired}
Let an attained positive maximum have $n\ge47$, $C_E\le C\le.4247$,
$0\le n(C-C_E)\le.7$, decrement at most $.19$, and both deleted
score caps at most $.4253$.
Given the shared contact support and strip implications, the exact
schedule above places its law in one of the 38 selected left rows or
45 selected central rows, with every resource in that row holding
simultaneously. This assertion has no upper bound on $n$.
\end{corollary}
\begin{proof}
The shared all-threshold support theorem supplies radius $3.84$.
Apply the finite enclosure theorem in the displayed order, retaining
the weighted bounds at every edge. The initialization's stored upper
index 98 occurs only through $.4095+K/98$; here $K=0$, so it imposes
no upper-index restriction. All subsequent index-dependent inequalities
are checked for every $n\ge47$ by their stated monotonicity gates.
The accepted closed covers exclude the complementary cells and give
the stated 83 rows.
\end{proof}

\paragraph{The local-comparison interface.}
For each final row supply $p\in I$, $\theta\le\Theta$, $\lambda\le\Lambda$,
the actual support/amplitude, $d_Y\le d_0$, both signed phase limits,
and all four coherent quadratics bounding \eqref{j:local:budget}.
On a phase/noise subcell, let $T_H$ be the maximum of each concave
quadratic at its clipped vertex, then the maximum over corners and the
minimum with the scalar bound. If a proved moment inequality is
\[
 D\ge\alpha d_Y+b\theta+P,\qquad\alpha>0,
 \quad \ell\ge\ell_-,\quad\lambda\ge L_0,
\]
then the exact strict resource supplied to the comparison is
\begin{equation}\label{j:local:re:localG}
 nd_Y<G:=\frac{T_H/\kappa-NP-b\ell_-^2L_0}{\alpha}.
\end{equation}
This uses $n\theta=\ell^2\lambda$ and $nP\ge NP$. If $G\le0$,
the entire closed subcell is excluded. Formula~\eqref{j:local:re:localG} is
the same resource as \eqref{j:local:budget}, not an additional allowance.

\paragraph{One confinement lemma with two concentration choices.}\label{j:local:confinement}
A proved amplitude is used as a support update, never as a premise of
its own proof. The following derivation applies to every acquired row,
with the same formulas for the central and left probability ranges.
\paragraph{Statement and hypotheses.}
Let $n\ge N=47$, $Y=B+U$, $\Pp(B=1)=p<1/2$ and $\E(U\mid B)=0$.
Put $q=1-p$, $s=\E U^2$, $\lambda=ns$ and $X=\ell(Y-p)$, where
$\E X^2=1$. Thus $\ell^2=(1-\theta)/(pq)$ and $\theta=\ell^2s$.
The conditional label-zero support is at most $a\ge0$ and the label-one
support is at least $b\le1$, where $a,b$ are support points and
$b-a=1-d$. In particular $0\le a\le d$, $1-d\le b\le1$.
Assume the initial strip $\operatorname{supp}Y\subset(-.55,1.55)$.
Write $\rho=\E|X|^3$, $\beta=\E[X|X|]$ and $\epsilon=\E|U|$.

Select an acquired closed probability panel $[l,h]$.
Its fields give upper bounds $\Theta,d_0,\rho_0,\beta_+,T,\Lambda$,
a lower bound $\beta_-$, and separate bounds $H_+,H_-$ for the two signs of
$H$. Assume
\[
 \theta\le\Theta,\quad d\le d_0,\quad \rho\le\rho_0,
 \quad \beta_-\le\beta\le\beta_+,\quad 0\le\lambda\le\Lambda,
 \quad .4<C\le C_0=.4253.
\]
For the positive sign $0\le H\le H_+$; for the negative sign
$-H_-\le H\le0$. Let $x=np+H/\ell$ and
$\nu=\mathcal L(x-\sum_{i<n}Y_i)$. Assume a contact majorant $Q$ for
$g(y)=\nu([y,\infty))$, with equality at each support point of $Y$,
and derivative
\begin{equation}\label{j:local:conf:derivative}
 \sqrt n(-Q'(y))=\ell\phi(H/\sqrt n)
 +\frac\ell n\{3C\beta+(H\phi(H/\sqrt n)+3C\rho)t-3Ct|t|\},
 \quad t=\ell(y-p).
\end{equation}
Both actual deleted convolutions satisfy, separately,
\begin{equation}\label{j:local:conf:concentration}
 Q_j(w)\le\frac{\ell\phi(0)w+2C_0\rho}{\sqrt j},
 \qquad j=n-1,n-2,
\end{equation}
for every interval of width $w\ge0$.

Here is the precise common resource assumption. Put
$\kappa=997/15000<\phi(0)/6$, $J=\int_{[a,b)}(y-a)\nu(dy)$,
and note that the central gap contributes $m\ell J=(1-\theta)J$ to $L$.
Thus $L\ge(1-\theta)J$; retain $e\ge0$ and
\begin{equation}\label{j:local:conf:budget}
 n e\rho+n\phi(0)D/6+n^{3/2}L<\mathcal B(H),\qquad
 \mathcal B(H)=\min\{T,\max_j(A_j+F_{j,\pm}|H|-Q_jH^2)\}.
\end{equation}
The four quadratics are the table's \texttt{quadratics} rows, using
\texttt{Qbudget}, not \texttt{Qraw}. The Gaussian phase contribution
has already been included in \texttt{Qbudget}.

For each panel let $E,F_\ell,m_0$ be respectively strict lower, upper,
and lower rational bounds for
\[
 \sqrt{(1-\Theta)/(h(1-h))},\qquad
 1/\sqrt{l(1-l)},\qquad \sqrt{(1-\Theta)l(1-l)}.
\]
The bounds use the grid $10^{-6}$, with an extra unit outward. Define
\begin{align*}
 c_+&=2(5.16)E^2(1-h-d_0)^2,&
 c_-&=2(7.16)E^2(l-d_0)^2,\\
 b_*&=\min\{3(4.16)E(1-h-d_0),3(6.16)E(l-d_0)\},\\
 P&=m_0E^2\pos{3.5\operatorname{dist}([l,h],[.41886,.41887])-2.041d_0}^2.
\end{align*}
The simultaneous moment inequality has the following usable form:
\begin{equation}\label{j:local:conf:moment}
 D\ge c_+M_++c_-M_-+b_*\theta+P,\quad
 d=(M_+/p+M_-/q)/\ell,\quad
 \epsilon\le2(M_++M_-)/\ell,\quad M_\pm\ge0.
\end{equation}
To derive it, the endpoint expansion of \eqref{j:core:gap-defect} gives
$D\ge m(\sqrt{\sigma+1}A-\sqrt{\sigma-1}B)^2
+2(\sigma+2)B^2M_++2(\sigma+4)A^2M_-
+3(\sigma+1)B\theta_++3(\sigma+3)A\theta_-$.
Insert $A\ge E(l-d_0)$, $B\ge E(1-h-d_0)$ and the rational lower
coefficients displayed above. For the square, write
$A=\ell(p-a)$, $B=\ell(1-p-(1-b))$, with
$a+(1-b)=d\le d_0$. Its unshifted coefficient is
$(\sqrt{\sigma+1}+\sqrt{\sigma-1})(p-p_*)$;
$\sqrt{\sigma+1}+\sqrt{\sigma-1}>3.5$ and
$\sqrt{\sigma+1}<2.041$ prove the stated lower square $P$.
For completeness, on the positive class put $Z=X-B\ge0$.
Then $T_+=2B^2M_++3B\E[Z^2;X>0]+\E[Z^3;X>0]$ and
$R_+=BM_++\E[Z^2;X>0]$. Consequently
$(\sigma-1)T_++6BR_+\ge2(\sigma+2)B^2M_+
+3(\sigma+1)B\theta_+$, since $\E[Z^2;X>0]\ge\theta_+$.
The reflected expansion gives the negative-class bound.
The centered absolute deviation on each class is at most twice its
mean distance from the inner endpoint, proving the bound on $\epsilon$.
All quantities refer to the same marginal.

\paragraph{Concentration from an antichain.}
\begin{lemma}[Strict-gap concentration]
Suppose $Y=U_0+B(1+U_1-U_0)$, where $B$ is Bernoulli$(p)$ independent
of the pair $(U_0,U_1)$, and $1+U_1-U_0\ge c>0$ almost surely.
For independent copies $Y_1,\ldots,Y_j$ and any interval $I$ of diameter
strictly less than $c$,
\[
 \Pp\left(\sum_{i=1}^jY_i\in I\right)
 \le \max_{0\le k\le j}\binom jk p^k(1-p)^{j-k}.
\]
The same bound holds after adding any independent random variable.
\end{lemma}
\begin{proof}
Condition on all noise pairs. The sum is a fixed constant plus
$\sum B_iw_i$, with every $w_i\ge c$. The family $\mathcal F$ of subsets
whose weighted sums lie in $I$ is an antichain: distinct comparable
subsets have sums differing by at least $c$, larger than the diameter.
A uniformly random permutation of $\{1,\ldots,j\}$ defines the maximal
chain of its initial subsets. The chain contains a fixed $k$-subset with
probability $1/\binom jk$, and meets an antichain at most once. Hence
\[
 \sum_{S\in\mathcal F}\binom j{|S|}^{-1}\le1.
\]
Multiplying each summand by its binomial level mass proves the bound
conditionally. Integration removes the conditioning. Conditioning on an
additional independent variable proves the last assertion. This argument
includes all atom collisions and any open or closed interval endpoints.
\end{proof}
For the actual marginal we may couple the two conditional noises
independently and then sample an independent label. Since $U_0\le a$ and
$U_1\ge b-1$, the step is at least $b-a=1-d$.
Strict diameter is necessary: a Bernoulli law on $\{0,c\}$ assigns mass
one to $[0,c]$, exceeding its largest atom when $0<p<1$.

\paragraph{A rational cap for every remaining index.}
Put $J_0=N-2=45$ and define
\[
 r_L=l-(1-l)/J_0,\qquad r_R=h+h/J_0,\qquad
 v_* =\min\{r_L(1-r_L),r_R(1-r_R)\}.
\]
The checker verifies $0<r_L\le r_R<1$. Adjacent binomial-mass ratios
show that any mode $k$ at index $j\ge J_0$ satisfies
$(j+1)p-1\le k\le(j+1)p$, hence $r=k/j\in[r_L,r_R]$.
In particular $1\le k\le j-1$. Finite Stirling bounds give
\[
 \sqrt j\binom jk p^k(1-p)^{j-k}
 \le\frac{\exp(1/(12j)-jD(r\Vert p))}{\sqrt{2\pi r(1-r)}}
 <\frac{.4}{\sqrt{v_*}(1-1/(12J_0))}.
\]
Here $D(r\Vert p)\ge0$ follows from Jensen's inequality for the logarithm;
$\exp x\le(1-x)^{-1}$ for $0\le x<1$ follows term by term from the
power series. The finite Stirling inequalities used are
$\sqrt{2\pi m}(m/e)^m<m!<\sqrt{2\pi m}(m/e)^m\exp(1/(12m))$.
Choose the strict $10^{-6}$ grid upper square root $M$ of
$.16/[v_*(1-1/(12J_0))^2]$. Thus, whenever $w<1-d$,
\begin{equation}\label{j:local:conf:antichainconcentration}
 Q_j(w)<M/\sqrt j,\qquad j\ge n-2\ge J_0.
\end{equation}
This covers the $n-1$ full summands of $\nu$ and the $n-2$ full summands
remaining after one mark or deletion. An independent remaining conditional
noise preserves the cap. We do not replace the latter index by $n-1$.


Either the score-derived cap or the strict-gap cap may be selected on
an interval where its hypotheses hold. The first confinement uses the
score cap; later central and left schedules use the strict-gap cap.
For the left panels replace the conditional cubic coefficient 24 below
by 28; in general any $c_\rho$ satisfying
$c_\rho^2[l(1-l)]^3>[3(1-l)+.55]^2$ is valid.
\paragraph{Signed deletion, including atoms.}
Let $m=n-1$, $Z=x-\sum_{i=1}^mY_i$, $\mu=x-mp$ and let $K=\sum B_i$.
For a fixed removed label put
$T_b(I)=\Pp(x-\sum_{i<m}Y_i-b-U_b\in I)$.
Exchangeability and $mp-K=Z-\mu+\sum U_i$ give the exact identities
\[
 \nu(I)-T_1(I)=\frac{\E[(Z-\mu+\sum U_i)1_I(Z)]}{mp},\qquad
 T_0(I)-\nu(I)=\frac{\E[(Z-\mu+\sum U_i)1_I(Z)]}{mq}.
\]
There is no absolute value around the sum of noises.
For a full removed pair $(Y_i,U_i)$, condition on the pair. The remaining
sum gives a function $f(Y_i)\in[0,Q_{m-1}(w)]$. Since $\E U_i=0$,
\[
 \E[U_i f(Y_i)]\le Q_{m-1}(w)\E(U_i)_+
 =\tfrac12\epsilon Q_{m-1}(w).
\]
This also works when one conditional law is deterministic.

Let independent conditional noises satisfy $\Delta=U_1-U_0\ge-d$.
For $\tau>0$ put $I^+=I+1+[-d,\tau]$. On $\Delta\le\tau$, the
point for $T_1$ in $I$ maps to the point for $T_0$ in $I^+$. On the
complement, conditioning on both noises gives error at most
$Q_{m-1}(w)\Pp(\Delta>\tau)$. Moreover
\begin{equation}\label{j:local:conf:bad}
 \Pp(\Delta>\tau)\le
 \min\left\{\frac{s}{s+p\tau^2},\frac d{d+\tau}\right\}.
\end{equation}
Indeed $\E\Delta=0$, $\operatorname{Var}\Delta\le s/p$, and
$\Delta+d\ge0$ has mean $d$. Cantelli and Markov prove the two bounds.
The exceptional event is strict, so $\Delta=\tau$ belongs to the good event.

If $I,I^+$ are disjoint, combine the two signed noise terms before
bounding them. For the same remaining-sum law their conditional
probabilities $f_I,f_J$ satisfy
\[
 0\le f_I\le Q_I,\quad0\le f_J\le Q_J,\quad f_I+f_J\le Q_{\rm hull}.
\]
Since $p<q$,
\[
 \frac{f_I}{p}+\frac{f_J}{q}
 \le\min\left\{\frac{Q_I}p+\frac{Q_J}q,
 \frac{Q_{\rm hull}}q+\left(\frac1p-\frac1q\right)Q_I\right\}.
\]
Multiplication by $\E U_+=\epsilon/2$ proves the combined bound.
Closed, open and half-open intervals are all covered by containment in
the corresponding closed interval. No density or atom separation is used.

\paragraph{Slope bounds and the two exterior intervals.}
Split $|H|$ into 16 closed equal cells $[h_l,h_h]$ and $\lambda$ into
24 closed equal cells $[s_l,s_h]$. On each box put
\[
 \begin{gathered}
 \rho_* =\min\{\rho_0,\rho_B(l)+24s_h/N\},\\
 b_- =\max\{\beta_-,1-2h-2(1-h)\Theta\},\qquad
 b_+ =\min\{\beta_+,1-2l+\Theta\}.
 \end{gathered}
\]
Here $\rho_B(p)=(1-2p(1-p))/\sqrt{p(1-p)}$. Write $s_i=\E U_i^2$. Conditional centering and the two sign supports give
\[
 \rho=\ell^3\{pq(1-2pq)+3pq(s_0+s_1)
                  +p\E U_1^3-q\E U_0^3\}.
\]
The initial strip and $d_0<.55$ imply $|U_i|<.55$.
Thus $\rho\le\rho_B+\ell^3(3q+.55)s\le\rho_B+24s$; the last
coefficient inequality is checked as
$24^2[l(1-l)]^3>[3(1-l)+.55]^2$.
Similarly $\beta=(1-\theta)(1-2p)+\ell^2(ps_1-qs_0)$, and
$|\ell^2(ps_1-qs_0)|\le\theta$. These identities prove the stated
outward bounds on $\rho$ and $\beta$.
Set $f=.4F_\ell$, $\varphi_-=.39886(1-h_h^2/(2N))>0$.
The exact upper bound for $H\phi(H/\sqrt n)$ used below is
\[
 v_+=\begin{cases}.4h_h&H\ge0,\\-\varphi_-h_l&H\le0.\end{cases}
\]
For all $t\le0$, completing the square in \eqref{j:local:conf:derivative} gives
\[
 \mathcal M_-=F_\ell\pos{(v_++3C_0\rho_*)^2/(12(.4))-3(.4)b_-},
 \qquad L_-=E\varphi_- -\mathcal M_-/N.
\]
A positive bound $B_*>\sqrt{(1-l)/l}$ gives valid central-gap lower
and positive-half-line upper slope bounds
\begin{align*}
 \mathcal M_+&=F_\ell\pos{.4h_hB_*+3C_0B_*(B_*-1)-1.2b_-},\\
 L_g&=E\varphi_--\max(\mathcal M_-,\mathcal M_+)/N,\\
 U_+&=f+\frac{F_\ell}{N}
 \left\{3C_0b_++\frac{(.4h_h+3C_0\rho_*)^2}{12(.4)}\right\}.
\end{align*}
They follow by completing a square on the positive half-line and using
$t(t-1)\le B_*(B_*-1)$ for $1\le t\le B_*$; on $0\le t\le1$ that
term has the favorable sign. Negative values are covered by $\mathcal M_-$.
All divisions and positivity requirements are explicit gates.

\paragraph{The sharper lower-ray slope.}
Take $\tau_-=5r_0/12$, $\eta=10^{-5}$ and
$I=[a-r,a-d-\tau_--\eta]$, where $r\ge r_0$.
Since $a\le d$, the entire interval has
$t\le t_*=-E(l+\tau_-+\eta)<0$.
Write $B(t)=b_-+\rho_*t+t^2$ and
$B_*^{\min}=B(\min(t_*,-\rho_*/2))$. For the actual $C$ and $t\le t_*$,
\[
 3C(\beta+\rho t+t^2)+H\phi t
 \ge1.2B(t)+v_+t+3(C_0-.4)\min(0,B_*^{\min}).
\]
The inequality uses $t<0$ and separates the possible negative part of
$B(t)$ before varying $C$. Minimize the right side at the clipped vertex
$\min(t_*,-(v_++1.2\rho_*)/2.4)$ and call its value $q_*$. Then
\[
 L=E\varphi_--F_\ell\max(0,-q_*)/N
\]
is a lower slope on the whole interval, for every $n\ge N$.
Positive corrections are discarded. Therefore contact gives
\[
 \sqrt n\nu(I)\ge L(r-d-\tau_--\eta).
\]
This one slope replaces the previous separate bounds on the long and
short parts of the interval.

\paragraph{Lower interval geometry.}
Its translate is $I^+=[a-r+1-d,b-\eta]$, of width $r-\eta$.
The initial strip implies $I^+\subset[a+g,b)$,
$g=1-.55-2d_0>0$, hence $\nu(I^+)\le J/g$.
The intervals are disjoint throughout $r\le.55+d_0$, since
$\inf I^+-\sup I=1-r+\tau_-+\eta>0$.
At $r=r_0$, concentration and directed distance bounds are
\begin{align*}
 Q_I&=f(r_0-\tau_--\eta)+2C_0\rho_*,&
 Q_J&=f(r_0-\eta)+2C_0\rho_*,\\
 Q_{\rm hull}&=f(1+r_0-\eta)+2C_0\rho_*,\\
 D_I&=\pos{v_H-l-\tau_--\eta},&
 D_J&=\pos{v_H+1-l-\eta},
\end{align*}
where $v_H=-h_l/F_\ell$ for $H\ge0$ and $v_H=h_h/E$ for $H\le0$.
These signs follow from $\mu=p+H/\ell$.

\paragraph{Upper interval geometry.}
For a support point $y\ge1+r_0$, take $\tau_+=7r_0/15$,
$I=[a,a+\tau_+]$ and $I^+=[b,a+1+2\tau_+]$.
The strict gates $d_0+2\tau_+<r_0$ and $1-d_0-\tau_+>0$ give
$I^+\subset[b,y)$ and disjointness. Contact, including the atom at $b$,
and the leakage definition give
\[
 \sqrt n\nu(I)\ge L_g(1-d)-\sqrt nJ/\tau_+,\qquad
 \sqrt n\nu(I^+)\le U_+(.55+d).
\]
Use
\begin{align*}
 Q_I&=f\tau_++2C_0\rho_*,&Q_J&=f(2\tau_++d_0)+2C_0\rho_*,\\
 Q_{\rm hull}&=f(1+2\tau_+)+2C_0\rho_*,\\
 D_I&=\pos{d_0+\tau_++h_h/E-l},&
 D_J&=\pos{1+d_0+2\tau_++h_h/E-l}.
\end{align*}
The upper side deliberately keeps the absolute phase bound.


\paragraph{The strict-gap alternative.}
If $1-.55-2d_0>0$ and the displayed lower intervals have width below
$1-d$ throughout $r\in[r_0,.55+d_0]$, take $Q_I=Q_J=M$.
For the upper intervals the corresponding gate is
$1-2d_0-2\tau_+>0$; then again $Q_I=Q_J=M$.
Use the same directed distances. Their concentration caps are now
independent of $r$, and use
$a_e=F_dM\{1/l+1/(1-h)\}/2$ in the next formulas.
The lower-ray gate reduces to $L>0$: its contact side increases with
$r$, while the entire error is constant on the acquired finite strip.
For the score-derived cap retain the explicit affine slope gate below.
No assertion outside that strip is needed.

\paragraph{One budget, two error bounds.}
Let $F_d$ be the strict grid upper bound for
$\sqrt{\max\{(47/46)^3,47/45\}}$. It simultaneously covers the two
concentration factors in \eqref{j:local:conf:concentration} after multiplication by
$\sqrt n$, uniformly for $n\ge47$. For either pair of intervals set
\begin{align*}
 a_e&=\frac{F_d}{2}\min\left\{\frac{Q_I}l+\frac{Q_J}{1-h},
 \frac{Q_{\rm hull}}{1-h}+\left(\frac1l-\frac1{1-l}\right)Q_I\right\},\\
 A&=F_d\left\{Q_ID_I/l+Q_JD_J/(1-h)\right\},\qquad B=F_dQ_I.
\end{align*}
The probability denominators are outward because $1/p-1/(1-p)$ is
decreasing on $(0,1/2)$.
Use $a_d=L$, $a_J=1/g$ on the lower side, and
$a_d=L_g+U_+$, $a_J=1/\tau_+$ on the upper side.
For the phase cell, maximize each concave budget quadratic at its
vertex clipped to $[h_l,h_h]$, take their maximum, then the minimum
with $T$; call this upper bound $T_*$. Put
\[
 R=T_*/\kappa-b_*E^2s_l,
 \quad\alpha=\min(c_+lE,c_-(1-h)E),\quad G=(R-NP)/\alpha.
\]
If $G\le0$, the strict common budget excludes the entire closed box.
Otherwise $nd<G$ and
\[
 B_\tau=\min\{s_h/(s_h+Nl\tau^2),G/(N\tau+G)\}
\]
bounds \eqref{j:local:conf:bad}. With
\[
 z=\max\left\{\frac{a_d/l+2a_e}{Ec_+},
 \frac{a_d/(1-h)+2a_e}{Ec_-},\frac{\kappa a_J}{1-\Theta}\right\},
\]
one valid total error bound is
\begin{equation}\label{j:local:conf:oldbound}
 \mathcal E_0=(zR+A)/N-zP+B B_\tau.
\end{equation}
This follows by multiplying the single budget by $z/\kappa$.
The inverse-index coefficient $zR+A$ is checked positive.

A stronger bound spends the same resource on the pair event. Let
$R_0=R-NP>0$ and define the three pairs
\[
 (v_+,w_+)=\left(\frac{a_d/l+2a_e}{Ec_+}\frac{R_0}{N},
                    \frac{R_0}{Ec_+l}\right),
\]
\[
 (v_-,w_-)=\left(\frac{a_d/(1-h)+2a_e}{Ec_-}\frac{R_0}{N},
                    \frac{R_0}{Ec_-(1-h)}\right),\qquad
 (v_J,w_J)=\left(\frac{\kappa a_J}{1-\Theta}\frac{R_0}{N},0\right).
\]
They are upper cost/displacement coordinates at the vertices of the
simplex with resources $nc_+M_+$, $nc_-M_-$ and $n^{3/2}L/\kappa$.
For any $t\ge0$,
\[
 \frac{D}{N\tau+D}\le\frac{N\tau}{(N\tau+t)^2}D+
 \frac{t^2}{(N\tau+t)^2}\qquad(D\ge0).
\]
The difference is $N\tau(D-t)^2/[(N\tau+D)(N\tau+t)^2]$.
Consequently a further valid error bound is
\[
 \mathcal E_t=\frac AN+
 \max_{i\in\{+,-,J\}}\left(v_i+B\frac{N\tau}{(N\tau+t)^2}w_i\right)
 +B\frac{t^2}{(N\tau+t)^2}.
\]
The checker takes the minimum of \eqref{j:local:conf:oldbound} and these eleven
bounds with $t=jG/10$, $j=0,\ldots,10$. Each is individually valid;
no numerical minimization or derivative of their minimum is assumed.

The lower contradiction is the strict gate
$Lr_0-L(\tau_-+\eta)>\mathcal E$.
To extend it to every $r\ge r_0$, bound increments directly in the original
affine-in-width signed deletion bound for the same law. The budget gives
$n\epsilon<E_*=2T/(\kappa E\min(c_+,c_-))$. The sufficient slope gate is
\[
 L>\frac{F_df}{N}\left\{D_I/l+D_J/(1-h)
 +\frac{E_*}{2}\left(1/l+1/(1-h)\right)\right\}+F_dfB_\tau.
\]
Using the uncoupled increment is safe even when the hull bound was used
at $r_0$. The upper contradiction is $L_g-.55U_+>\mathcal E$.
Thus no support point lies below $a-r_0$ or above $1+r_0$.
Since $a\ge0$, $b\ge1-d_0$ and $d_0<r_0$, the label-zero and label-one
supports both satisfy $|U|<r_0$.


\paragraph{The selected schedules.}
The ordinary, central-antichain, and left-antichain certificates substitute
their recorded $r_0$ into these same gates. Noise-dependent amplitudes
are the same theorem with a positive lower noise endpoint in
$R=T_*/\kappa-b_*E^2s_l$. In the central rows the additional amplitudes
are $.35,.30,.275$ on $[.05,.08],[.08,.10],[.10,\Lambda]$;
on the left they are $.375,.325,.30$. Intersect each interval with its
own row's $\Lambda$ and each amplitude with that row's global amplitude.
The fixed acceptance checks both signs of $H$, all closed phase/noise
cells, clipped endpoints, and strict budget exclusions. No proof of
LOW, MID, or a binomial score comparison is used by this confinement lemma.




\paragraph{Finite-index substitution in confinement.}
\label{j:local:finite-confinement}
Here the numerical constants in the confinement proof of
Section~\ref{j:local:confinement} are parameters. More precisely,
replace its floor $47$ by $N\ge39$, its initial strip $.55$ by a
number $w\in(0,1)$, and its conditional cubic coefficient by
$c_\rho>0$. Retain a probability panel $[l,h]\subset(0,1/2)$ and assume
\[
 \operatorname{supp}Y\subset(-w,1+w),\quad d_0<w,
 \quad c_\rho^2[l(1-l)]^3>[3(1-l)+w]^2.
\]
The strip and the inward bounds $U_0\le d_0$, $U_1\ge-d_0$
together give $|U_i|<w$. The displayed conditional cubic identity gives
$\rho\le\rho_B+\ell^3(3q+w)s\le\rho_B+c_\rho s$.
Use $J_0=N-2$ for the binomial mode interval and the deletion factor
$\sqrt{\max\{(N/(N-1))^3,N/(N-2)\}}$. In the interval geometry
use $g=1-w-2d_0>0$ and $r\le w+d_0$; in the upper contradiction
use $L_g-wU_+>\mathcal E$. The other width, disjointness,
positive-slope and budget gates are unchanged. These substitutions
follow directly through the proof: $w$ enters only the conditional
cubic bound and the endpoints of the known strip, while the index
enters only the decreasing factors already bounded above. They
therefore prove the same confinement implication whenever all the
resulting gates pass.

The positive-allowance calculation at actual index $39$ uses $N=39$,
$J_0=37$ and deletion factor
$\sqrt{\max\{(39/38)^3,39/37\}}$. Its acquired initial widths are
$w=.55$ or $.575$; the central and left cubic coefficients are
$24$ and $30$, respectively. All the same explicit gates are checked
with these parameters before a contraction is used.
The central floor-$40$ sequence uses $N=40$, $w=.55$ initially and
$c_\rho=24$. The left rows start with their separately acquired
strip $w\in\{.55,.575,.6\}$ and use $c_\rho=30$. After a strict
contraction, the proved amplitude becomes the next strip width.
Thus every later use has an already established strip, and the
left sequence does not assume the narrower uniform-tail strip.



\subsection{The common comparison identity and the nearest count}\label{j:local:shortloc:section}

A law concentrated near two points has a natural integer count and a small
centered error. We compare it with its own Bernoulli count. The same exact
identity handles every noise regime. There are two useful estimates of its
right side: discard all but the nearest outgoing count, or retain all counts
through a random cutoff. The estimates share their normalization, moments,
and tail bounds; only the final scalar witness changes.

\paragraph{The identity and the common numerical test.}

Let $0<p\le1/2$, $q=1-p$, $v=pq$, and let
\[
 Y=B+U,\qquad \Pr(B=1)=p,\qquad \mathbb E(U\mid B)=0.
\]
The two actual conditional laws $U_0,U_1$ obey
\begin{equation}\label{j:local:shortloc:support}
 |U_j|\le A,\quad U_0\le a_0,\quad U_1\ge-a_1,\quad
 a_0,a_1\ge0,\quad d=a_0+a_1\le d_0<\min(p,q),\quad d_0\le A.
\end{equation}
Write $s_j=\mathbb E U_j^2$, $s=qs_0+ps_1$, $\lambda=ns$,
$\ell=(v+s)^{-1/2}$, $\theta=s/(v+s)$ and
$\rho=\mathbb E|\ell(Y-p)|^3$. At the actual threshold choose
\[
 x=np+H/\ell=k+u,\qquad k\in\{1,\ldots,n-1\},\quad |u|\le\tfrac12,
 \qquad \delta=k-np=H/\ell-u.
\]
Either nearest integer is allowed at a tie. Put
\[
 \pi_j=\binom njp^jq^{n-j},\quad F_B(j)=\Pr(\operatorname{Bin}(n,p)\le j),
 \qquad \rho_B=(1-2v)/\sqrt v,
\]
with $\pi_j=0$ outside $[0,n]$. Let $W_j$ be an independent sum of $j$
copies of $U_1$ and $n-j$ copies of $U_0$. The signed scores are
\[
 C=\frac{\sqrt n}{\rho}\{\Pr(\textstyle\sum Y_i\le x)-\Phi(H/\sqrt n)\},
 \qquad C_B=\frac{\sqrt n}{\rho_B}\{F_B(k)-\Phi(\delta/\sqrt{nv})\}.
\]
Throughout this section $C\ge c_*>0$; the selected application uses $c_*=.4$.
Define
\begin{align}
 f(u)&=\frac{\Phi((\delta+u)/\sqrt{n(v+s)})-
                    \Phi(\delta/\sqrt{n(v+s)})}{\pi_k},\label{j:local:shortloc:f}\\
 Q&=\Pr(\textstyle\sum Y_i\le k+u)-F_B(k),\nonumber\\
 \mathcal R&=\sqrt n\{\Phi(\delta/\sqrt{nv})-
                 \Phi(\delta/\sqrt{n(v+s)})\}+C(\rho_B-\rho).\nonumber
\end{align}
Direct subtraction gives the exact identity
\begin{equation}\label{j:local:shortloc:master}
 \rho_B(C-C_B)=\sqrt n\pi_k\{Q/\pi_k-f(u)\}+\mathcal R.
\end{equation}
In particular, the phase $u$, the sign of $H$, and both conditional laws
refer to the same marginal throughout.

\begin{samepage}
\begin{theorem}[Parameterized local comparison]\label{j:local:shortloc:theorem}
Fix $n\ge N$, $\lambda>0$ and a closed parameter cell on which
\eqref{j:local:shortloc:support} holds. Suppose $\sqrt n\pi_k\ge m>0$.
Either of the following witnesses proves $\rho_B(C-C_B)<-\varepsilon\lambda$.
\begin{enumerate}
\item \emph{Nearest-count witness.} Valid bounds $O,T,R_0,\kappa$ satisfy
\[
 \Pr(W_k>u)+f(u)\ge O,\quad
 \sum_{i\ge1}\Pr(W_{k+i}\le u-i)\le T,\quad
 \pi_{k+i}/\pi_k\le R_0,\quad \mathcal R\le\kappa\lambda/n.
\]
Only legal counts occur in the sum; adding positive upper bounds for
other counts is allowed. If $D=O-R_0T>0$ and
\begin{equation}\label{j:local:shortloc:directtest}
             mD-\kappa\lambda/n>\varepsilon\lambda,
\end{equation}
the conclusion follows.
\item \emph{Whole-count witness.} Bounds $B(\lambda)<0$ and $I(\lambda)\ge0$
satisfy
\[
 \rho_B(C-C_B)\le\sqrt n\pi_k\lambda B(\lambda)
                                      +\lambda I(\lambda)/n.
\]
It is enough that
\begin{equation}\label{j:local:shortloc:whole-test}
                  mB(\lambda)+I(\lambda)/N<-\varepsilon.
\end{equation}
\end{enumerate}
The witnesses may differ from cell to cell. A finite closed cover and a
separate treatment of $\lambda=0$ therefore suffice for a complete local
comparison.
\end{theorem}
\end{samepage}
\begin{proof}
Conditioning on the binomial count gives, with all atom conventions intact,
\begin{equation}\label{j:local:shortloc:count}
 Q=-\sum_{i=0}^{k}\pi_{k-i}\Pr(W_{k-i}>u+i)
       +\sum_{i=1}^{n-k}\pi_{k+i}\Pr(W_{k+i}\le u-i).
\end{equation}
Dropping the outgoing terms with $i>0$ gives $Q/\pi_k-f(u)\le-D$.
Use \eqref{j:local:shortloc:master} and $D>0$. For the second witness, use
$B<0$ before replacing $\sqrt n\pi_k$ by $m$, and $I\ge0$ before replacing
$n$ by $N$. These sign requirements are part of the test, not conventions.
\end{proof}

\paragraph{Normalization and the repair, once for every witness.}

Here is an explicit finite recipe for the constants required above.
Let $p\in[p_-,p_+]\subset(0,1/2)$, $n\ge N$, $\lambda\le\Lambda$,
$\theta\le\Theta$, and $|\delta|\le\Delta$, where $\Delta\ge1$ also includes
every binomial mode. Choose $0<r_-<p_- -\Delta/N$ and
$p_++\Delta/N<r_+<1/2$. Let $v_*>0$ be a lower bound for $p(1-p)$.
Choose positive $E,L,R$ bounding, on $[r_-,r_+]$,
\[
 |D'''(r\Vert p)|=\frac{1-2r}{r^2(1-r)^2},\qquad
 |(\log(r(1-r)))'|=\frac{1-2r}{r(1-r)},\qquad
 \frac1{12r(1-r)},
\]
respectively, where $D(r\Vert p)=r\log(r/p)+(1-r)\log((1-r)/q)$.
Then, with $Z(t)=Lt/2+R+Et^3/(6N)$,
\begin{equation}\label{j:local:shortloc:mass}
 \sqrt n\pi_k=\frac{\phi(0)}{\sqrt v}
       e^{-\delta^2/(2nv)+\zeta},\qquad n|\zeta|<Z(|\delta|).
\end{equation}
The finite Stirling estimate
$\log(j!)=(j+1/2)\log j-j+\log(2\pi)/2+r_j$, $0<r_j<1/(12j)$,
used in the concentration argument above gives this formula on substitution
in $n!/(k!(n-k)!)$ and Taylor expansion of relative entropy and the logarithmic
prefactor. Here is a finite proof of the estimate. For
$a_j=\log(j!)-(j+1/2)\log j+j$ and $z=(2j+1)^{-1}$,
\[
0<a_j-a_{j+1}=\sum_{i\ge1}\frac{z^{2i}}{2i+1}<\frac1{12j(j+1)}.
\]
Summation bounds $a_j-\lim a_i$. The integration-by-parts recursion for
$\int_0^{\pi/2}\sin^j t\,dt$, and the inequalities between successive
integrals, give Wallis's product and $\lim a_i=\log(2\pi)/2$.
The entropy Taylor error and factorial error give respectively
$E|\delta|^3/(6N)$ and $1/[12r(1-r)]$ in $Z$.

Consequently a valid mass lower bound is any positive
\begin{equation}\label{j:local:shortloc:m}
 m<\frac{.39886}{\sqrt{p_+(1-p_+)}}
          \left[1-\frac{\Delta^2/(2v_*)+Z(\Delta)}N\right].
\end{equation}
The same formula at a mode proves $\sqrt n\max_j\pi_j<1$ whenever
$v_*[1-Z(\Delta)/N]^2>.4^2$ and $Z(\Delta)<N$.
The rational bounds $.39886<\phi(0)<.4$ follow, for example, by the
alternating arctangent series in Machin's identity.

Differentiating \eqref{j:local:shortloc:f} gives
\begin{equation}\label{j:local:shortloc:logder}
 \log f'=-\tfrac12\log(1+s/v)+\delta^2/(2nv)
                   -(\delta+u)^2/(2n(v+s))-\zeta.
\end{equation}
Thus $n|\log f'|\le B_0\lambda+P(D,a)$ on $|\delta|\le D\le\Delta$,
$|u|\le a\le1/2$, where
\[
 B_0>\frac1{2v_*}+\frac{\Delta^2}{2Nv_*^2},\qquad
 P(D,a)=\frac{2Da+a^2}{2v_*}+Z(D).
\]
Choose $0<b<1$ above $[B_0\Lambda+P(\Delta,1/2)]/N$ and put
$F=(1-b)^{-1}$. The elementary estimates $e^z\le(1-z)^{-1}$ and
$1-e^{-z}\le z$ imply
\begin{equation}\label{j:local:shortloc:derivative}
 |f'-1|\le(B_1\lambda+P_1)/n,\qquad B_1=FB_0,\quad P_1=FP(D,a).
\end{equation}
This supplies both a derivative interval $l_f\le f'\le h_f$ and the
finer phase-error bound $|u-f(u)|\le a(B_1\lambda+P_1)/n$.

The exact count variance is
\begin{equation}\label{j:local:shortloc:variance}
 t_j=\mathbb E W_j^2=\lambda+(j-np)(s_1-s_0),\qquad
 |s_1-s_0|\le s/p\le c_v s,\quad c_v\ge1/p_-.
\end{equation}
It gives $\omega_-\lambda\le t_k\le\omega_+\lambda$ with
$\omega_\pm=1\pm c_vD/N$. All uses below require the checked condition
$\omega_->0$. Also
$t_{k+i}\le iV_*\lambda$ for $V_*=1+c_v(\Delta+1)/N$.
The binomial recursion gives, for every legal $i\ge1$,
\begin{equation}\label{j:local:shortloc:ratio}
 \frac{\pi_{k+i}}{\pi_k}\le R_0,
 \qquad R_0>1+\frac{p_++\Delta-1}
 {Nv_*[1-(\Delta-1)/(Np_-)]},
\end{equation}
provided the denominator is positive and $p_++\Delta<2$.
Indeed the first ratio has this bound and all subsequent rising factors
are at most one. No count outside $[0,n]$ is used in this recursion.

Finally \eqref{j:local:shortloc:support} fixes the signs inside the absolute cubic.
Expansion and centering give
\[
 \mathbb E|Y-p|^3=v(1-2v)+3pq(s_0+s_1)
                   +p\mathbb E U_1^3-q\mathbb E U_0^3
 \ge v(1-2v)+(3p-d)s.
\]
The inequality $(1+z)^{3/2}\le1+3z/2+3z^2/8$ for $z\ge0$ yields
\begin{equation}\label{j:local:cubic-bound}
 (v+s)^{3/2}(\rho-\rho_B)\ge s
     \left(6p-3p^2-\tfrac32-d-\frac{3(1-2v)s}{8v}\right).
\end{equation}
Let $Q_*>0$ be a lower bound for the bracket and
$m_c\le c_*/[p_+(1-p_+)+\Lambda/N]^{3/2}$.
The normal variance correction is nonpositive for $\delta\le0$;
for $\delta>0$ its derivative in $s$ is at most $.4\delta/(2v^{3/2})$.
Therefore
\begin{equation}\label{j:local:shortloc:repair}
 \mathcal R\le(P_*D_+-m_cQ_*)\lambda/n,
 \qquad P_*\ge\frac{.4}{2v_*^{3/2}},\quad D_+\ge\delta_+.
\end{equation}
One may take $D_+=D$, but retaining the positive part is often decisive.
All bounds concern $v+s$, which can exceed $1/4$.

\paragraph{Two elementary tools for probabilities.}

\begin{lemma}[A shifted quartic]\label{j:local:shortloc:quartic}
Suppose $W$ is a sum of independent centered variables bounded in absolute
value by $A$, and $t=\mathbb E W^2$. Let $K_0\ge(9+6\sqrt3)/4$ and $a>0$.
For $r\ge0$ and $\sigma\in\{-1,1\}$,
\begin{equation}\label{j:local:shortloc:quarticformula}
 \Pr(W+\sigma r>0)\ge
 \frac{(3a^2-A^2)t-3t^2-4Art-6r^2t-r^4
                      +3a^2r^2+2\sigma a^3r}{K_0a^4}.
\end{equation}
\end{lemma}
\begin{proof}
The polynomial $h_a(z)=(-z^4+3a^2z^2+2a^3z)/(K_0a^4)$ is below
$\mathbf1_{\{z>0\}}$: it is nonpositive for $z\le0$, and its maximum for
$z>0$ is at $z/a=(1+\sqrt3)/2$ and is at most one. Independence gives
$|\mathbb E W^3|\le At$ and $\mathbb E W^4\le A^2t+3t^2$.
Expand $\mathbb Eh_a(W+\sigma r)$ and apply these two inequalities.
The value at the threshold is zero, so atoms require no limiting argument.
\end{proof}

\begin{lemma}[One common tail bound]\label{j:local:shortloc:tails}
If independent centered $X_i\le A$ have total variance at most $V$, then,
for $b>0$, their sum has upper tail bounded by the minimum of one and
\begin{equation}\label{j:local:shortloc:tailmin}
 \begin{gathered}
 \frac{V}{V+b^2},\qquad
 \mathcal B(A,V,b)=e^{b/A}\left(\frac{V}{V+Ab}\right)^{b/A+V/A^2},\qquad
 \frac{\mathbb E(P_\mu-j)_+^3}{(b/A+\mu-j)^3},\\
 \mu=V/A^2,\qquad 0\le j<b/A+\mu.
 \end{gathered}
\end{equation}
where $P_\mu$ is Poisson. The same holds for a lower tail when $|X_i|\le A$.
All bounds apply to closed events, hence also to strict events.
\end{lemma}
\begin{proof}
Cantelli's bound follows by applying Markov to $(\sum X_i+V/b)^2$.
Taylor's integral remainder gives
$\mathbb E e^{zX_i}\le\exp\{\mathbb E X_i^2(e^{zA}-1-zA)/A^2\}$
for $z\ge0$. Multiply and use $z=A^{-1}\log(1+Ab/V)$ to get $\mathcal B$.
The zero-variance case is immediate.

For the third bound, smooth $f(x)=(x-t)_+^3$ by a nonnegative smooth
compactly supported density. Then $f'',f''',f''''\ge0$. For
$\mathcal Af(x)=f(x+A)-f(x)-Af'(x)$, the Taylor quotient
\[
 \frac{f(x+y)-f(x)-yf'(x)}{y^2}
       =\int_0^1(1-r)f''(x+ry)\,dr
\]
is nondecreasing in $y$. A centered $X\le A$ of variance $v$ consequently
satisfies $\mathbb E f(x+X)\le f(x)+(v/A^2)\mathcal Af(x)$.
For $\mathcal P_zf(x)=\mathbb E f(x+A(P_z-z))$,
$\partial_z\mathcal P_zf=\mathcal P_z\mathcal Af$ and
\[
 \partial_z^2\mathcal P_zf=\mathcal P_z\mathcal A^2f\ge0,\qquad
 \mathcal A^2f(x)=\int_0^A\!\int_0^A(A-r)(A-w)f''''(x+r+w)\,dr\,dw.
\]
Thus $\mathcal P_{v/A^2}f\ge f+(v/A^2)\mathcal Af$. Replace the independent
variables successively; translations and expectations preserve the derivative
signs and Poisson rates add. Increasing the total variance to $V$ is permitted
by Jensen after adding independent centered Poisson noise. Removing the
smoothing is justified by bounded actual upper supports and finite Poisson
third moments. Markov at $t=A(j-\mu)$ proves the displayed bound.
\end{proof}

These bounds have finite rational evaluations. For example,
\begin{align}\label{j:local:shortloc:poissonfinite}
 \mathbb E(P_\mu-j)_+^3={}&\mu^3+3(1-j)\mu^2+(1-3j+3j^2)\mu-j^3\\
 &+e^{-\mu}\sum_{h=0}^{\lfloor j\rfloor}\frac{\mu^h(j-h)^3}{h!}.\nonumber
\end{align}
Replacing $e^{-\mu}$ by $(\sum_{h=0}^{24}\mu^h/h!)^{-1}$ enlarges the
nonnegative correction. Alternatively sum the positive Poisson series
through 32; the remaining tail is at most its first term divided by
$1-\mu(34-j)^3/[34(33-j)^3]$, provided $j<33$ and that ratio is below one.
For Bennett, put $x=Ab/V$ and
$L=[(b/A+\mu)\log(1+x)-b/A]$. The positive series
$\log(1+x)=2\sum_{h\ge0}z^{2h+1}/(2h+1)$, $z=x/(2+x)$, gives lower
bounds for $L$; any positive lower exponential series gives an upper bound
for $e^{-L}$. Rounding a positive lower $L$ down before exponentiation is
also valid. These are choices of rational witnesses, not optimization claims.

\paragraph{The nearest-count witness, including vanishing noise.}

For $u=-r\le0$, add $-h_fr$ to \eqref{j:local:shortloc:quarticformula} with
$\sigma=1$; for $u=r\ge0$, add $l_fr$ with $\sigma=-1$.
Both bound $\Pr(W_k>u)+f(u)$ from below. For a cell
$t\in[t_-,t_+]$, $r\in[r_-,r_+]$, this expression is concave in $t$
(second derivative $-6/(K_0a^4)$), so only $t_-,t_+$ are needed.
At each endpoint, bound a monomial with positive coefficient at $r_-$
and one with negative coefficient at $r_+$. Any rational $a>0$ is a valid
witness. Other valid choices are
\begin{equation}\label{j:local:shortloc:ownother}
 \frac{r_-^2}{t_++r_-^2}-h_fr_+\quad(u\le0),\qquad l_fr_-\quad(u\ge0).
\end{equation}
They follow from Cantelli and nonnegativity, respectively. Their maximum
or one selected witness is valid without an optimality proof.

For incoming counts use $t_{k+i}\le iV_*\lambda$. On $u\in[-r_+,-r_-]$
the first distance is $1+r_-$ and every later distance is at least $i$.
On $u\in[r_-,r_+]$ they are $1-r_+$ and $i(1-r_+/2)$ for $i\ge2$.
Let $b,b_o$ be these first and later distances. On $\lambda\in[L,U]$,
let $T_1$ be any bound from Lemma~\ref{j:local:shortloc:tails} at $(A,V_*U,b)$.
Since $\mathcal B(A,iV,ib)=\mathcal B(A,V,b)^i$, all later counts total at most
$B^2/(1-B)$ for $B=\mathcal B(A,V_*U,b_o)<1$. Thus
\begin{equation}\label{j:local:shortloc:midtest}
 D=O-R_0\{T_1+B^2/(1-B)\}>0,\qquad
 mD-\kappa U/N>\varepsilon U
\end{equation}
proves the entire cell when $\kappa\ge0$. The selected MID normalizers
explicitly verify this sign. This is the common central and left MID proof.

A cell touching $\lambda=0$ requires estimates proportional to $\lambda$;
a positive lower endpoint may not be substituted. The following rescaling
provides them. Put $r=zt$, $t_*=\omega_+U$ and $d_a=K_0a^4$.
The shifted quartic plus $br$, divided by $t$, becomes
\begin{equation}\label{j:local:shortloc:scaledown}
 F_{\sigma,b}(z,t)=\frac{3a^2-A^2-3t}{d_a}
 +\left(\frac{2\sigma}{K_0a}+b-\frac{4At}{d_a}\right)z
 +\frac{3a^2t-6t^2}{d_a}z^2-\frac{t^3z^4}{d_a}.
\end{equation}
It extends to $t=0$ and is concave in $t$. On $0\le z\le z_0$, replacing
$z^4$ by $z_0^2z^2$ gives a lower quadratic at $t=t_*$; at $t=0$ it is
linear. Endpoints and a clipped quadratic vertex therefore give an exact
lower bound $J(\sigma,b,z_0)$ for the whole rectangle.

For negative phase choose $z_-,c_->0$ and verify
\begin{gather}\label{j:local:shortloc:tinynegative}
 J(1,-h_f,z_-)>c_-,\quad
 \frac{z_-^2}{1+z_-^2t_*}-h_fz_->c_-,\quad
 \frac{2z_-}{(1+z_-^2t_*)^2}>h_f,\\
 \frac1{1+4t_*}-\frac{h_f}{2}>c_-t_*,\qquad z_-t_*<1/2.\nonumber
\end{gather}
The quartic handles $r\le z_-t$. Beyond that point use
$r^2/(t+r^2)-h_fr$. Its derivative at $r=z_-t$ is positive by the third
gate; its second derivative has just one sign change, so any later critical
point is a maximum. Its minimum is at $z_-t$ or $1/2$, covered by the other
gates. Hence the own count exceeds $c_-t$.
For positive phase choose $z_+,c_+>0$ with
\begin{equation}\label{j:local:shortloc:tinypositive}
 J(-1,l_f,z_+)>c_+,\qquad l_fz_+>c_+,\qquad z_+t_*<1/2.
\end{equation}
The quartic handles $r\le z_+t$ and the normal term handles the rest.
The own count exceeds $c_+t$ and also exceeds $l_fr$.

For incoming tails proportional to $\lambda$, fix a Poisson anchor $j$.
A positive series term of index $h$ divided by
$\lambda(b/A+\mu-j)^3$ has logarithmic derivative, multiplied by $\lambda$,
\[
 h-1-\mu-\frac{3\mu}{b/A+\mu-j},\qquad \mu=V_*\lambda/A^2.
\]
If $b/A-j>0$ and this is nonnegative at $\lambda=U$ for the first positive
index, it is nonnegative for every term and all $0<\lambda\le U$.
Therefore the Poisson bound divided by $\lambda$ is largest at $U$.
Likewise $B^2/[\lambda(1-B)]$ is increasing whenever
$\lambda B'/B\ge1/2$. Differentiating the Bennett formula gives
\[
 \lambda B'/B=b/A-\mu\log(1+b/(A\mu));
\]
The function $g(\mu)=b/A-\mu\log(1+b/(A\mu))$ decreases because
$\log(1+x)>x/(1+x)$ for $x>0$. Moreover the logarithmic derivative of
$B^2/[\lambda(1-B)]$, multiplied by $\lambda$, is
$2g-1+gB/(1-B)\ge0$ when $g\ge1/2$. Thus checking the gate at $U$
proves the whole interval. A rational upper bound for the logarithm
verifies that gate. On cells using the fixed positive bound $l_fr_-$ instead,
unscaled endpoint tails suffice: after a positive bracket is checked, divide
it by $U$ to obtain a lower bound per $\lambda$. These arguments, followed by
\eqref{j:local:shortloc:directtest}, prove TINY uniformly down to arbitrary positive
noise, including a deterministic conditional component.

\subsection{Whole-count comparison}

Assume $c=1-d_0>1/2$. Couple all $2n$ conditional variables independently and
set $W_j=\sum_{i\le j}U_{1,i}+\sum_{i>j}U_{0,i}$, extending it constantly
outside $[0,n]$. Every increment of $j+W_j$ is at least $c$. Define
\[
 K=\max\{j\in\mathbb Z:j+W_j\le k+u\},\qquad R=K-k.
\]
An independent binomial selector has the correct actual count laws, so
\begin{equation}\label{j:local:shortloc:cutoff}
                  \Pr(\textstyle\sum Y_i\le k+u)=\mathbb EF_B(k+R).
\end{equation}
Let $\chi=(n+1)p-k<2$ and $m_B=\max_j\pi_j$. For every integer $r$,
\begin{equation}\label{j:local:shortloc:curvature}
 F_B(k+r)-F_B(k)-r\pi_k\le\frac{m_B}{(n+1)v}
 \left[(\chi-1)_+r_++\chi_+\binom{r_-}{2}+\binom{r_-}{3}\right].
\end{equation}
Indeed, the $\operatorname{Bin}(n+1,p)$ mass is
$\pi_j^+=q\pi_j+p\pi_{j-1}\le m_B$. Factorial cancellation gives
\[
 \pi_j-\pi_{j-1}=\frac{(n+1)p-j}{(n+1)v}\,\pi_j^+
 \qquad(j\in\mathbb Z),
\]
including the exterior zero masses. Thus the positive part at $j=k+i$
is at most $m_B(\chi-i)_+/((n+1)v)$.
Double summation gives \eqref{j:local:shortloc:curvature}: for $r>0$ only $i=1$
can contribute; for $r=-h<0$ use $(\chi+i)_+\le\chi_++i$ and sum the
multiplicities $h-1-i$. Zero extension makes the same calculation valid
outside the legal count range.

Write $t=\mathbb EW_k^2$. Crossing monotonicity and Cauchy--Schwarz give
\begin{equation}\label{j:local:shortloc:crossingmoments}
 \mathbb ER_+\le\frac{t}{(c-u)^2},\quad
 \mathbb E\binom{R_-}{2}\le\frac{t}{[c+\min(u,c/3)]^2},\quad
 \mathbb E\binom{R_-}{3}\le\frac{t\sqrt{A^2+3t}}{(2c-1/2)^3}.
\end{equation}
For the first, $R=i\ge1$ implies $-W_k\ge ci-u$ and
$(ci-u)^2\ge i(c-u)^2$. For $R=-h\le-2$, the next count is outside the
closed event, so $W_k>c(h-1)+u$. Put $j=h-1$.
The function $j(j+1)/(2(cj+u)^2)$ decreases for $j\ge1$ if $u\le c/3$;
for larger $u$ reduce to that endpoint. Also
$(j^3-j)/(6(cj-1/2)^3)$ decreases for $j\ge2$, since its derivative has
the sign of $2cj-(3j^2-1)/2<0$. These give the last two bounds using
$\mathbb E(W_k)_+^3\le\sqrt{\mathbb EW_k^2\mathbb EW_k^4}$.

\paragraph{The marked crossing.}
Let $D_i=U_{1,i}-U_{0,i}$, $\vartheta=\mathbb E(D_i)_+=\mathbb E(D_i)_-$,
$0<\eta\le1$, and $R_f=\mathbb E\{u+\eta-W_k\}$, where
$\{x\}=x-\lfloor x\rfloor$ is the fractional part. The two deleted laws are
\[
 A_0\overset{d}=kU_1+(n-k-1)U_0,\qquad
 A_1\overset{d}=(k-1)U_1+(n-k)U_0.
\]
Then
\begin{align}
 \mathbb ER&\le u+\eta-R_f+J,\label{j:local:shortloc:mark}\\
 J&=\frac{\vartheta}{\eta}\sum_{i\ge1}
 \{\Pr(A_0\le u+d_0-ci)+\Pr(A_1>u+\eta+i-A)\}.\nonumber
\end{align}
Here are the pathwise details. Put $z=u-W_k$ and mark
\[
 Z_f=\sum_{i=1}^{n-k}(D_{k+i})_-\mathbf1_{\{W_k\le u-ci\}},\quad
 Z_b=\sum_{i=1}^{k}(D_{k-i+1})_+\mathbf1_{\{W_k>u+\eta+i\}}.
\]
If $R>z+\eta$, the crossing activates the marks through $R$: for $R>0$
use $z\ge cR$, while for $R<0$ use $z<R-\eta$. In either case
$R\le z+Z_f+Z_b$. Hence if $h=R-\lfloor z+\eta\rfloor>0$,
$Z_f+Z_b>h-1+\eta\ge\eta h$. Thus
$R\le\lfloor z+\eta\rfloor+(Z_f+Z_b)/\eta$ on every path.
On an active forward mark $U_0>U_1\ge-d_0$; deleting its pair implies
$A_0\le u+d_0-ci$. On a backward mark delete $U_1\le A$ to obtain the
strict event for $A_1$. The remaining sum is independent of the entire
pair, so the expectations factor. Extending to all $i\ge1$ adds only
nonnegative bounds. Finally $|R|\le1+(|W_k|+|u|)/c$ proves integrability
and permits expectations. No coupling has replaced either deleted marginal.

\paragraph{A finite bound for every mark.}
Fix $\lambda\in[L,U]$, $L>0$, and $t\le\omega_+\lambda$; put
$T_L=\omega_+L$, $T_U=\omega_+U$. On a phase subcell $u\in[v_l,v_h]$,
for $i=1,2$ let $b_1=ci-d_0-v_h>0$ and $b_2=i-A+v_l>0$.
Lemma~\ref{j:local:shortloc:tails} gives a bound $\omega_+\lambda K_i^{\rm ind}$,
where
\[
 K_i^{\rm ind}=\sum_{j=1}^2\min\left\{\frac1{b_j^2+T_L},
                          \frac{\overline P(A,T_U,b_j)}{T_L}\right\}.
\]
The denominator $T_L$ is a lower bound for the chosen upper envelope
$\omega_+\lambda$, not for any actual variance.
Alternatively couple $A_0=Z+U_1$, $A_1=Z+U_0$ with
$Z\overset d=(k-1)U_1+(n-k-1)U_0$. The two events project into
$Z\le-\alpha$ and $Z>\beta$, where
$\alpha=ci-2d_0-v_h>0$, $\beta=i-A-d_0+v_l>0$.
For any $z\in(-\alpha,\beta)$ their sum is at most
$(\omega_+\lambda+z^2)/\min(z+\alpha,\beta-z)^2$.
Consequently a second kernel is
\[
 K_i^{\rm com}=\frac{1+z^2/T_L}{\min(z+\alpha,\beta-z)^2},\qquad
 z=\max\{-T_L/\beta,\min((\beta-\alpha)/2,T_L/\alpha)\}.
\]
This center is inside the interval; differentiating the two quadratic
ratios on either side of the midpoint proves its optimality, although
any interior rational center suffices. Take the smaller kernel on each
phase subcell and the maximum over their closed cover. Denote it $K_i$.
For $i\ge3$ use fourth moments and the convex midpoint integral for
$(cx-b)^{-4}$. On $u\in[u_l,u_h]$ this gives
\begin{equation}\label{j:local:shortloc:K}
 K=K_1+K_2+\frac{A^2+3T_U}{3c}
  \left\{\frac1{((5/2)c-A-u_h)^3}
                 +\frac1{((5/2)c-A+u_l)^3}\right\},\qquad
 J\le\frac{\vartheta}{\eta}\omega_+\lambda K.
\end{equation}
All denominators must be positive. Indeed, the tail fourth moment is at
most $A^2t+3t^2$, and integration from $5/2$ to infinity covers every
remaining integer, including the harmless added illegal counts.

If the same-law resource supplies $nd<G$ and
$G^2<c_vNL$, then
\begin{equation}\label{j:local:shortloc:theta}
 \vartheta\le\frac{c_vG\lambda}{c_vn\lambda+G^2}.
\end{equation}
To prove this, for centered $X\ge-b$, $\mathbb EX^2\le V$ and $b^2\le V$,
the quadratic $b(X-V/b)^2/(b+V/b)^2$ majorizes $X_-$ on $[-b,\infty)$.
Its expectation is at most $bV/(V+b^2)$. Apply this to $D_i\ge-d$,
using $\operatorname{Var}(D_i)=s_0+s_1\le c_v\lambda/n$, and monotonicity
in $b\le\sqrt V$. The case $d=0$ is immediate. With $\eta=z\lambda$,
$z>0$, $zU\le1$, \eqref{j:local:shortloc:K} becomes
\begin{equation}\label{j:local:shortloc:Jfinal}
 \frac J\lambda\le\frac{\omega_+Kc_vG}{z(c_vN\lambda+G^2)}.
\end{equation}

\paragraph{The fractional-part gain.}
If each centered summand has support span at most $b<1/2$ and the sum has
variance $t$, then, for any real $w$ and
$0<r_*\le2\sin^2(\pi b)/b^2$,
\begin{equation}\label{j:local:shortloc:harmonic}
             \mathbb E\{w-W\}\ge\frac5{23}(1-e^{-r_*t}).
\end{equation}
Indeed $1-\cos(2\pi x)\le(23/5)x$ for $0\le x\le1$.
A finite proof of this sole scalar inequality is recorded below.
For an independent copy $Z'$ of one summand, concavity of sine on
$[0,\pi b]$ gives
\[
 |\mathbb Ee^{2\pi iZ}|^2
 =1-2\mathbb E\sin^2\pi(Z-Z')
 \le e^{-4\sin^2(\pi b)\operatorname{Var}(Z)/b^2}.
\]
Multiply square roots over the independent summands and use periodicity.
This argument includes zero variances and integer collisions.
Here the conditional spans are at most $A+d_0$, so this is the $b$ to use.

For a reproducible scalar proof, $x\ge10/23$ follows from $1-\cos\le2$.
On $[0,10/23]$ use the degree-ten alternating upper polynomial for
$1-\cos(2\pi x)$, with directed coefficients:
\[
 Q(x)=\sum_{j=1}^5\frac{(-1)^{j+1}(2\pi_jx)^{2j}}{(2j)!},\qquad
 \pi_j=\begin{cases}3.141592653589794,&j\text{ odd},\\
                       3.141592653589793,&j\text{ even}.\end{cases}
\]
Put
$P(x)=23/5-Q(x)/x$. On the eight intervals with left endpoint $a_j=5j/92$
and width $w=5/92$, its degree-nine Bernstein coefficients are
\[
 b_{jk}=\sum_{i=0}^k\frac{\binom ki}{\binom9i}
           \sum_{h=i}^9[x^h]P(x)\binom hi a_j^{h-i}w^i,
 \qquad 0\le j<8,\quad0\le k\le9.
\]
The 80 rational coefficients are positive; the retained scalar acceptor
reconstructs them. Expansion in powers of the unit-interval coordinate
proves the formula; Bernstein weights are nonnegative and sum to one.
The alternating remainder and the stated rational bounds on $\pi$ follow
from the same elementary series above. This completely specifies the
finite scalar obligation. A rational rate follows by replacing
$\sin(\pi b)$ with $x-x^3/6+x^5/120-x^7/5040$, $x=3.14159b$, on the
checked alternating-series domain.

Set $r=r_*\omega_-$, require $rU<3$, and define
\begin{equation}\label{j:local:shortloc:gain}
 g=\frac5{23}\left(\frac{r^3L^2}{6}-\frac{r^4L^3}{24}
                     +\frac{r^5L^4}{120}-\frac{r^6L^5}{720}\right).
\end{equation}
The sixth-order Taylor lower bound for $1-e^{-r\lambda}$ gives
$R_f/\lambda\ge5r/23-5r^2\lambda/46+g$. The discarded higher-order
part is increasing as a function of $\lambda$ for $r\lambda<3$:
its derivative has the sign of $48-18x+(24/5)x^2-x^3>0$ on $[0,3]$.
For instance this polynomial has positive degree-three Bernstein
coefficients $48,30,132/5,51/5$ there. Thus freezing it at $L$ is valid.

\paragraph{One final expression.}
On a signed phase cell find an interval $[x_l,x_h]$ containing $H/\ell$
by its four endpoint quotients, using positive bounds for $\ell$.
Let
\[
 D=\max(|x_l-u_h|,|x_h-u_l|),\quad
 D_+=\max(0,x_h-u_l),\quad a=\max(|u_l|,|u_h|),\quad b_\chi=p_+-x_l.
\]
The same formula allows coarser upper bounds $D,D_+,b_\chi$, as used by
the central certificate. Check $D\le\Delta$ and $b_\chi+u_h<2$.
With $\omega_\pm=1\pm c_vD/N$, the curvature coefficient is
\begin{equation}\label{j:local:shortloc:V}
 V_0=\frac{\omega_+}{v_*}\left[
 \sup_{u_l\le u\le u_h}\left\{
 \frac{(b_\chi+u-1)_+}{(c-u)^2}
 +\frac{(b_\chi+u)_+}{[c+\min(u,c/3)]^2}\right\}
 +\frac{\sqrt{A^2+3T_U}}{(2c-1/2)^3}\right].
\end{equation}
On a rational phase subcell, both numerators are bounded at the right
endpoint; the first denominator is smallest there, and the second is
smallest at the left endpoint. This proves the supremum bound, including
both kinks. Equations \eqref{j:local:shortloc:cutoff}--\eqref{j:local:shortloc:crossingmoments},
$\sqrt n m_B\le1$, and \eqref{j:local:shortloc:master} now give the whole-count
witness with
\begin{align}\label{j:local:shortloc:BI}
 B(\lambda)&=z+\frac{\omega_+Kc_vG}{z(c_vN\lambda+G^2)}
                      -\frac{5r}{23}+\frac{5r^2\lambda}{46}-g,\\
 I(\lambda)&=V_0+aB_1+\frac{aP_1}{\lambda}+P_*D_+-m_cQ_*.
 \nonumber
\end{align}
We used $|u-f(u)|\le a(B_1\lambda+P_1)/n$ and $\sqrt n\pi_k\le1$.
Every constant is fixed before $\lambda$ varies. Both $B$ and
$mB+I/N$ are convex on $\lambda>0$; $I$ is decreasing. Thus their two
endpoint gates $B<0$, $I>0$, $mB+I/N<-.001$ prove the entire closed noise
cell. This proves all three LOW variants by exactly the same argument.

\subsection{Signed refinements and the complete local covers}

The retained tables contain 45 central probability intervals covering
$[.38,.47]$ and 38 left intervals covering $[.304,.38]$. Each row supplies
its own $\Theta$, signed $H$ bounds, $d_0$, $\Lambda$, and global amplitude;
the confinement profiles may improve the amplitude with $\lambda$.
All these are hypotheses of the present section, acquired earlier without
using its conclusion. The shared resource has the form
\[
 nD_{\rm mom}<\frac{15000}{997}T_H,\qquad
 D_{\rm mom}\ge\alpha d+\beta\theta+P.
\]
Consequently a noise cell $[L,U]$ has
\[
 nd<G=\{(15000/997)T_H-NP-\beta\ell_-^2L\}/\alpha.
\]
Here $T_H$ is the maximum of the acquired concave quadratic budget on the
closed signed $H$ cell, further bounded by its scalar ceiling. It already
contains the phase charge. If $G\le0$, the entire cell is impossible and
receives no score test; the exclusion is strict even at $G=0$.

For positive noise the complete cover is the following. Each interval is
intersected with the row ceiling $\Lambda$; shared endpoints are included.
\begin{center}\small
\begin{tabular}{@{}llll@{}}\toprule
Noise & Central phase & Left phase & Witness\\\midrule
$(0,.05]$ & all & all & nearest count (scaled)\\
$[.05,.06]$ & $|u|\ge.09$ & $|u|\ge.09$ & nearest count\\
$[.06,.08]$ & $|u|\ge.13$ & $|u|\ge.14$ & nearest count\\
$[.08,.10]$ & $|u|\ge.18$ & $|u|\ge.19$ & nearest count\\
$[.10,.12]$ & $|u|\ge.29$ & $u\le-.29$ & nearest count\\
$[.10,.12]$ & --- & $u\in[.29,.5]$ & whole count\\
$[.05,\Lambda]$ & remaining phase & remaining phase & whole count\\\bottomrule
\end{tabular}
\end{center}
Beyond $.12$ the last row covers the full phase interval. The three LOW
objects are one central cover, one left cover with phase slices, and one
extra positive left cover. Their distinction is only the closed domain
and the chosen outward constants. Signed $H$ cells meet at zero; nearest
count ties meet at $u=\pm1/2$. The fixed coverage checker also checks every
row-ceiling singleton and every entirely excluded LOW row.

The common TINY parameters are
\[
 N=47,\quad p\in[.304,.47],\quad |H|\le1.393,\quad
 \Theta=.0154,\quad d_0=.0181,\quad A=.425,\quad\Lambda=.05.
\]
For its scaled own witness use $a=.4784$, $K_0=9/(4\cdot.4641)$,
$(z_-,c_-)=(1.66,.514)$ and $(z_+,c_+)=(1.2,1.036)$.
Its six phase cells are $[-.5,0],[0,.1],\ldots,[.4,.5]$.
The first two use scaled Poisson anchors $1.25,1$, respectively; the other
four use the fixed positive normal bound. The exact normalizers and gates
specified in the accompanying certificate map verify these constants.
The central MID amplitudes are $.35,.35,.30,.275$ in its four rows.
The left ones are $.375,.375,.325,.30$, each further reduced by the row's
global amplitude. No interpolation of the tables is used.

All six accepted comparisons give the same strict conclusion
\[
                        \rho_B(C-C_B)<-.001\lambda.
\]
Their proofs have now reduced to Theorem~\ref{j:local:shortloc:theorem}, the shared
normalization, and the two explicitly derived witnesses. The arithmetic
objects still have their original domains and targets; no certificate
branch has been silently deleted or relabeled.

If $\lambda=0$, conditional centering and $p,q>0$ imply $U_0=U_1=0$
almost surely. The law is exactly Bernoulli. Schulz's strict finite-binomial
bound~\cite[Theorem~1]{Schulz2016} states, for every $0<p<1$, finite $n$
and real $x$,
\[
 \sqrt n\left|\Pp(\operatorname{Bin}(n,p)\le x)
 -\Phi((x-np)/\sqrt{npq})\right|<\CE\rho_B.
\]
It applies directly to the actual threshold, including atoms. At positive
noise the same theorem closes $C<C_B<\CE$. No estimate divides by zero
or replaces a uniform small-noise bound with a fixed-law limit.

\paragraph{A triangle of allowed noise and displacement.}
Introduce Bernoulli coordinates
\[
 X=\ell(B_0+U-p),\quad B_0=\mathbf1_{\{X>0\}},\quad
 \ell=\bar A+\bar B=\frac m{pq},\quad \E[U\mid B_0]=0.
\]
Then
\[
 \ell^2=\frac{1-\theta}{pq},\quad m\ell=1-\theta,
 \qquad d=\frac{\Delta_A+\Delta_B}{\ell},\quad
 \lambda=n\E U^2=\frac{npq\theta}{1-\theta}.
\]

\begin{theorem}[Joint inward/noise constraint]\label{j:local:triangle}
Suppose $n\ge N>0$, and bounds for this same law give
\[
 nD\le N d_D,\quad \theta\le\Theta<1,\quad pq\le v_*,
 \quad A+B\ge S_*>0,
\]
and
\[
 D\ge b+6mS_*(\Delta_A+\Delta_B)+K\theta,
 \qquad b\ge0,\quad K>0.
\]
Set $r=d_D-b$. If $r>0$, define
\[
 d_* =\frac{r}{6(1-\Theta)S_*},\qquad
 \Lambda=\frac{Nv_*r}{K(1-\Theta)}.
\]
Then
\begin{equation}\label{j:local:eq:triangle}
                 \boxed{\displaystyle
                   \frac d{d_*}+\frac\lambda\Lambda\le1.}
\end{equation}
If $r<0$, the premises are impossible. If $r=0$, they force
$d=\lambda=0$.
\end{theorem}
\begin{proof}
Multiply the assumed joint lower bound for $D$ by $n$. Since all
coefficients being weakened are nonnegative,
\begin{align*}
 Nd_D\ge nD
 &\ge nb+6n(1-\theta)S_*d+K\frac{(1-\theta)\lambda}{pq}\\
 &\ge Nb+6N(1-\Theta)S_*d+\frac{K(1-\Theta)}{v_*}\lambda.
\end{align*}
Subtract $Nb$. For $r>0$, division by $Nr$ proves
\eqref{j:local:eq:triangle}. For $r\le0$, positivity of both coefficients gives
the stated alternatives without dividing by zero.
\end{proof}

Apply this with the binary charge in \eqref{j:core:simultaneous},
$S_*=A_l+B_l$ and $K=\min\{3(\sigma+1)B_l,3(\sigma+3)A_l\}$.
\begin{corollary}[Closed noise bands and outward rounding]
If $0\le L\le\lambda\le U\le\Lambda$, then
\[
                       d\le d_*(1-L/\Lambda).
\]
Both intercepts in \eqref{j:local:eq:triangle} may first be replaced by larger
positive numbers. The same conclusion then holds using those enlarged
intercepts and the enlarged noise ceiling.
\end{corollary}
\begin{proof}
Rearrange \eqref{j:local:eq:triangle} and use $\lambda\ge L$.
Increasing either denominator weakens its nonnegative summand, so it
preserves the joint inequality. This also proves the rounding assertion.
\end{proof}



\paragraph{The coefficient must be bounded in the correct direction.}
Suppose $n\ge N$, $p\in[p_-,p_+]\subset(0,1/2)$,
$0\le d\le d_0$, and $0\le\lambda\le\Lambda$.
Choose $0<v_*\le p_-(1-p_-)$ and put
\[
 V^*=p_+(1-p_+)+\Lambda/N,\qquad
 Q_*=6p_--3p_-^2-\frac32-d_0
                     -\frac{3(1-2v_*)\Lambda}{8v_*N}.
\]
Then $Q\ge Q_*$: the first term increases with $p$, and
$(1-2v)/v$ decreases with $v$. In particular, $Q_*$ may be negative.
Already with $d=s=0$ its sign changes at $p=1-1/\sqrt2$.

Let $C$ be the score multiplying the moment correction, and assume
$0<c_*\le C\le c^*$. Choose positive multipliers satisfying
\[
        0<m_-\le\frac{c_*}{(V^*)^{3/2}},\qquad
        m_+\ge\frac{c^*}{v_*^{3/2}},
\]
and define the signed gain
\begin{equation}\label{j:local:gain}
             G_*:=\min\{m_-Q_*,m_+Q_*\}.
\end{equation}

\begin{proposition}[Signed gain and nonnegative repair]\label{j:local:repair}
Let $\Phi$ denote the standard normal distribution function. For an actual
count displacement $\delta$, set
\[
 \mathcal R=\sqrt n\left\{
 \Phi\!\left(\frac{\delta}{\sqrt{nv}}\right)
 -\Phi\!\left(\frac{\delta}{\sqrt{n(v+s)}}\right)\right\}
                         +C(\rho_B-\rho).
\]
If $D_+\ge\max(\delta,0)$ and $P_*\ge .4/(2v_*^{3/2})$, then
\begin{equation}\label{j:local:repair_bound}
 \mathcal R\le(P_*D_+-G_*)\frac{\lambda}{n}
 \le \kappa\frac{\lambda}{N},\qquad
 \kappa=\max\{0,P_*D_+-G_*\}.
\end{equation}
\end{proposition}
\begin{proof}
Write $A=C/(v+s)^{3/2}$. The hypotheses give
$0<m_-\le A\le m_+$. By \eqref{j:local:cubic-bound},
$C(\rho-\rho_B)\ge sAQ\ge sAQ_*$.
If $Q_*\ge0$, $AQ_*\ge m_-Q_*$; if $Q_*<0$,
$AQ_*\ge m_+Q_*$. This proves $C(\rho-\rho_B)\ge sG_*$
in both cases. One must use the \emph{upper} multiplier in the negative
case.

The normal part is nonpositive when $\delta\le0$.
For $\delta>0$, differentiation with respect to an intermediate variance
$t\in[0,s]$ gives
\[
 \frac{d}{dt}\left[-\sqrt n\,
 \Phi\!\left(\frac{\delta}{\sqrt{n(v+t)}}\right)\right]
 =\frac{\delta}{2(v+t)^{3/2}}
   \phi\!\left(\frac{\delta}{\sqrt{n(v+t)}}\right)
 \le\frac{.4\delta}{2v_*^{3/2}}\le P_*D_+.
\]
Integrate over $[0,s]$, and combine with the moment bound.
Finally $s=\lambda/n$ and $\kappa\ge0$, so
$(P_*D_+-G_*)\lambda/n\le\kappa\lambda/n\le\kappa\lambda/N$.
The clamp at zero is needed before replacing $n$ by its lower bound.
\end{proof}



\paragraph{A binomial mass bound on the actual phase cell.}
Suppose the finite Stirling argument has supplied positive constants
$E,L_0,R$ on a fixed segment containing every $k/n$ under consideration, with
\[
 \sqrt n\,\pi_k=\frac{\phi(0)}{\sqrt v}
       \exp\{-\delta^2/(2nv)+\zeta\},\qquad
 n|\zeta|\le Z_N(|\delta|),\qquad
 Z_N(t)=\frac{L_0t}{2}+R+\frac{Et^3}{6N},
 \quad n\ge N.
\]
The constants may have been computed on a much wider global phase range.
Let $v_-\le v\le v_+$ and let $q_+\ge\sqrt{v_+}$.
\begin{lemma}[Mass localization]\label{j:local:lem:mass}
On any smaller closed cell where $|\delta|\le D$,
\[
 \sqrt n\,\pi_k\ge
 \frac{.39886}{q_+}
 \left[1-\frac{D^2/(2v_-)+Z_N(D)}N\right].
\]
The maximum of this expression and any previously proved mass lower bound
is also a valid mass lower bound. There is no requirement $D\ge1$.
\end{lemma}
\begin{proof}
Set $A_*=D^2/(2v_-)+Z_N(D)\ge0$. The logarithmic formula gives
$\sqrt n\pi_k\ge(.39886/q_+)e^{-A_*/n}$.
Now $e^{-z}\ge1-z$ and $n\ge N$. If the bracket is negative the displayed
inequality still holds, and the previously positive bound can be retained.
The global Stirling segment, rather than the smaller value $D$, continues
to justify $E,L_0,R$.
\end{proof}

\paragraph{A sharper elementary harmonic constant.}
\begin{lemma}\label{j:local:lem:cos}
For every $x\ge0$,
\[
               1-\cos(2\pi x)\le\frac{4553}{1000}x.
\]
\end{lemma}
\begin{proof}
Put $K=4553/1000$. For $x\ge2/K$ the assertion follows from
$1-\cos\le2$. On $[0,2/K]$ take the degree-fourteen alternating upper
Taylor polynomial $Q(x)$ for $1-\cos(2\pi x)$. Its remainder is nonpositive:
the first omitted term is negative, and the remaining term magnitudes
decrease, since $(4\pi/K)^2<16\cdot15$.

Use the directed endpoints
\[
 \pi_-=3.141592653589793<\pi<3.141592653589794=\pi_+.
\]
In each positive coefficient of $Q$ replace $\pi$ by $\pi_+$; in each
negative coefficient replace it by $\pi_-$. This only increases $Q$.
The exact checker expands $P(x)=K-Q(x)/x$, a polynomial of degree thirteen,
on each of the 64 intervals $[2j/(64K),2(j+1)/(64K)]$. If
$P(x)=\sum_{h=0}^{13}p_hx^h$, its Bernstein coefficients on $[a,a+w]$ are
\[
 b_k=\sum_{i=0}^k\frac{\binom ki}{\binom{13}i}
        \sum_{h=i}^{13}p_h\binom hi a^{h-i}w^i,
               \qquad k=0,\ldots,13.
\]
All 896 rational coefficients are positive; the smallest is greater than
$.00012312$. Thus $P>0$ on the full interval. The value at zero follows by
continuity (or directly).

For completeness the stated $\pi$ endpoints are checked using alternating
series for $\arctan(1/5)$ (32 terms) and $\arctan(1/239)$ (12 terms).
The tangent addition formula gives
$\tan(4\arctan(1/5)-\arctan(1/239))=1$.
The angle lies between zero and $4/5<\pi/2$, so it equals $\pi/4$.
The positivity follows already from
$4(1/5-1/(3\cdot5^3))-1/239>0$; the upper bound follows from
$\arctan z<z$. This proves Machin's identity on the correct branch.
The full rational recipe is \texttt{check\_harmonic4553.py}; no sampled
positivity is used.
\end{proof}

If the original characteristic-function argument gives
$|\E e^{2\pi iW}|\le e^{-rt}$, Lemma~\ref{j:local:lem:cos} upgrades its fractional
part estimate from $(5/23)(1-e^{-rt})$ to
\[
                 \frac{1000}{4553}(1-e^{-rt}).
\]
Consequently every occurrence of $5/23$ and $5/46$ in that harmonic gain is
replaced by $1000/4553$ and $500/4553$, respectively. The existing degree-six
Taylor gain and its monotonicity argument remain unchanged.



\paragraph{The signed Gaussian correction in nearest-count comparison.}

Let $n\ge N$, $v=p(1-p)>0$, $s\ge0$, $\lambda=ns$, and let the
actual nearest count and phase satisfy
\[
 \delta=k-np=h-u,\qquad
 f(t)=\frac{\Phi((\delta+t)/\sqrt{n(v+s)})-
                   \Phi(\delta/\sqrt{n(v+s)})}{\pi_k}.
\]
The actual $k$ and $\delta$ are held fixed when $t$ varies. Assume the
proved finite Stirling formula on this count range is
\[
 \sqrt n\,\pi_k=\frac{\phi(0)}{\sqrt v}
       \exp\{-\delta^2/(2nv)+\zeta\},\qquad
 n|\zeta|\le Z_N(D),\quad |\delta|\le D.
\]
Write $v\ge v_->0$ and $Z=Z_N(D)\ge0$.

\begin{lemma}[Negative phase]
Suppose $u=-r$, $0\le a\le r\le b$, and $h\le h_*$. Define
\[
 t_*=\max(0,\min(b,h_*+b)),\qquad G_-=2(h_*+b)t_*-t_*^2,\qquad
 \eta_-=G_-/(2v_-)+Z.
\]
If $D^2\le Nv_-$ and $\eta_-<N$, then
\[
 f'(-t)\le\frac1{1-\eta_-/N}\quad(0\le t\le r),\qquad
 u-f(u)\le\frac{b\eta_-}{n(1-\eta_-/N)}.
\]
\end{lemma}
\begin{proof}
Direct differentiation and the Stirling formula give
\[
 \log f'(-t)=-\tfrac12\log(1+s/v)
 +\frac{\delta^2s}{2nv(v+s)}
 +\frac{2\delta t-t^2}{2n(v+s)}-\zeta.
\]
Since $\log(1+x)\ge x/(1+x)$ for $x\ge0$, the first two terms are at most
\[
 \frac{s}{2(v+s)}\left(\frac{\delta^2}{nv}-1\right)\le0.
\]
This cancellation removes the variance term from the upper derivative
bound. Because $\delta=h+r$, $r\le b$ and $0\le t\le r\le b$,
\[
 2\delta t-t^2\le 2(h_*+b)t-t^2\le G_-.
\]
The last maximum is the concave quadratic maximum on $[0,b]$, achieved
at the displayed projection $t_*$. It is nonnegative because $t=0$ is
available. Therefore $\log f'(-t)\le\eta_-/n$.
Using $e^x\le(1-x)^{-1}$ for $0\le x<1$ gives
\[
 f'(-t)-1\le\frac{\eta_-/n}{1-\eta_-/n}
             \le\frac{\eta_-}{n(1-\eta_-/N)}.
\]
Integrate $u-f(u)=\int_0^r(f'(-t)-1)\,dt$.
\end{proof}

Equivalently, $G_-=b^2+2h_*b$ for $h_*\ge0$, and
$G_-=\max(0,h_*+b)^2$ for $h_*<0$. This sharper vertex is used in the
selected finite left LOW implementation.

\begin{lemma}[Positive phase]
Suppose $u\in[a,b]\subset[0,1/2]$, $h\le h_*$, and define
\[
 G_+=\max\left(0,\max_{z\in[a,b]}(2h_*z-z^2)\right).
\]
Then
\[
 u-f(u)\le\frac b n
       \left(\frac\lambda{2v_-}+\frac{G_+}{2v_-}+Z\right).
\]
The inner maximum is at the projection of $h_*$ onto $[a,b]$.
\end{lemma}
\begin{proof}
The exact derivative identity and $\log(1+s/v)\le s/v$ give
\[
 \log f'(t)\ge-\frac{\lambda+2\delta t+t^2}{2nv}-\zeta.
\]
For fixed actual $u$, the function $2\delta t+t^2$ is convex on $[0,u]$.
Its maximum is at an endpoint, and $\delta=h-u$ gives
\[
 2\delta t+t^2\le\max(0,2hu-u^2)\le G_+.
\]
Use $e^{-x}\ge1-x$ and integrate from zero to $u$.
\end{proof}

\paragraph{Whole-count insertion and interval coverage.}
The whole-count identity is split as
$Q/\pi_k-f(u)=(Q/\pi_k-u)+(u-f(u))$.
Thus one needs an upper bound on the signed correction $u-f(u)$;
an absolute-value estimate is unnecessary. After division by $\lambda$,
the new coefficients replacing the old Gaussian correction are
\[
 \begin{array}{ll}
 u\ge0:&\displaystyle
 \frac b{2v_-}+\frac{b\{G_+/(2v_-)+Z\}}\lambda,\\[6pt]
 u\le0:&\displaystyle
 \frac{b\eta_-}{(1-\eta_-/N)\lambda}.
 \end{array}
\]
On each fixed phase cell these have the form $A+B/\lambda$ with
$A,B\ge0$. They therefore preserve the original convex endpoint check on
each closed noise interval. Every phase cell is confined to one half-line;
zero is a shared closed endpoint. The count displacement bound $D$ comes
from the same actual $h-u$ cell, and the Stirling constants continue to
come from the already proved global count segment.

The other all-index steps are unchanged: the event bracket is checked
negative before substituting a binomial mass lower bound, and the entire
inverse-$n$ correction is checked positive before replacing $n$ by $N$.
The inverse correction is $I(\lambda)=K+A+B/\lambda$, hence decreases with $\lambda$; positivity at both endpoints (in particular the upper endpoint) licenses the all-index step throughout the interval. The final score is convex in $\lambda$: its event fraction is a positive constant divided by an affine positive function, its $B/\lambda$ coefficient is nonnegative, and its remaining terms are affine or constant. These are independent gates in the exact kernel. The formulas apply to
both signs of $h$ without reflecting the law or changing the count during
integration. They also supply an improved upper derivative for the
negative-phase MID inequalities.
\paragraph{Full-phase MID below the whole-count threshold.}
For $\lambda\in[.05,.08]$, the selected shifted-quartic MID comparisons
are checked on every closed phase interval of length $.001$ in
$[0,1/2]$, for both phase signs. The shifted-quartic polynomial and
nearest-count identity are defined at phase zero, so the two closed covers
include $u=0$ directly. This permits the whole-count LOW comparison to
start at $\lambda=.08$. The four MID bands use phase starts
$0,0,.19,.27$ at noise bands $[.05,.06]$, $[.06,.08]$,
$[.08,.10]$, and $[.10,.12]$, respectively.




\paragraph{Positive allowances and finite inverse-index corrections.}
\label{j:local:allowance}
For actual index $39$ put $K=741/4000$. The target is
$C< C_B+K/39$, so the scalar tests differ from
the zero-allowance tests in Section~\ref{j:local:shortloc:section}.
All of them follow from the same identity
\eqref{j:local:shortloc:master}. Write $b=\sqrt{39}\pi_k$ and retain
the verified two-sided bound $m\le b\le1$. A nearest-count
bracket $D$ may now have either sign; consequently
\[
 bD\ge\min\{mD,D\}.
\]
If the repair is at most $r\lambda/39$, with $r\ge0$, a noise
band $[a,b_\lambda]$ is excluded by the intermediate-noise test
\[
 \min\{mD,D\}-r b_\lambda/39+K/39>0.
\]
Indeed the identity gives $\rho_B(C-C_B)<K/39$.
The elementary inequality $\rho_B\ge1$ follows from
$(1-2v)^2-v=(1-v)(1-4v)\ge0$, $0<v\le1/4$.
Thus $C<C_B+K/39\le\CE+K/39$.

At arbitrarily small positive noise the rescaled quartic and tail
estimates give the same inequality per unit $\lambda$. The tiny-noise
test adds the allowance
\[
             \frac{K}{39\max\{.05,\Lambda\}}
\]
to that per-noise reserve. Multiplying by
$0<\lambda\le\min\{.05,\Lambda\}$ again gives a bound
strictly below $K/39$. No positive lower noise cutoff is inserted.

For the whole-count test, write
\[
 S(\lambda)=\max\{mB(\lambda),B(\lambda)\}
                         +I(\lambda)/39.
\]
Both the bracket and the inverse correction may have either sign;
the first maximum uses $m\le b\le1$, while the division by $39$
uses the actual integer. Hence
$\rho_B(C-C_B)\le\lambda S(\lambda)$.
The saved gate on $[a,b_\lambda]$ is
\[
        K\rho_{B,l}/39-b_\lambda S(\lambda)>0,
\]
where $\rho_{B,l}$ is a positive lower bound for $\rho_B$.
If $S<0$ the conclusion is immediate; otherwise it follows by
$\lambda\le b_\lambda$ and $\rho_B\ge\rho_{B,l}$.
The formula is convex in positive $\lambda$, since its only
nonaffine terms are positive multiples of
$(\alpha\lambda+\beta)^{-1}$ and $1/\lambda$, and maxima of
convex functions preserve convexity. Both endpoint gates therefore
cover the whole noise cell.



At zero noise, the strict finite-binomial theorem gives $C<\CE$,
so every positive allowance is excluded there as well.
For a bounded integer band $N\le n\le L$, a signed inverse correction
has the valid bound $I/n\le\max(I/N,I/L)$. If
$b=\sqrt n\pi_k\in[b_0,1]$, a signed bracket satisfies
$bJ\le\max(b_0J,J)$ and $bJ\ge\min(b_0J,J)$.
These two-sided rules replace the sign restrictions needed by the
infinite-index test. The actual index remains in every binomial slack
calculation. With a uniform binomial reserve $\epsilon_p>0$ and an
upper bound $S$ on $\rho_B(C-C_B)/\lambda$, it is enough to verify
$\rho_{B,l}\epsilon_p>U\max(S,0)$ on a noise cell $\lambda\le U$.
No negative inverse correction is enlarged by replacing its denominator
with the lower index.

\paragraph{The local exclusion, including equality.}\label{j:local:comparison}
\begin{proposition}[Finite local exclusion]\label{j:local:finite42}
For an integer $42\le n\le48$, no attained positive ratio maximum satisfies
\begin{gather*}
 0<n(C_n-\CE)\le54/100,\qquad n(C_{n-1}-C_n)_+\le19/100,\\
 C_n<4226/10000,\qquad C_{n-1},C_{n-2}<4253/10000.
\end{gather*}
\end{proposition}
\begin{proof}
Apply the support polynomial with $(N,L,K,q,R)=(42,48,0,.19,3.84)$.
In particular its favorable fifth-order term is evaluated at $48$; the
infinite-range support row cannot be substituted for this finite band.
The subsequent bounded-threshold tests use excess allowance .54, loss .19
and own-score cap .4226. The resource invariant is
$\mathcal R_n\le r$, $nD\le42d$. The transfer argument of
Section~\ref{j:local:transfer} applies at floor $42$, including the signed
curvature, retained constant penalties and actual-noise triangle.
Confinement uses the already acquired strips and the substitutions in
Section~\ref{j:local:finite-confinement}, with deletion floor $40$.
The left radius first contracts to $3.025$; contact then proves the
coordinate strip $(-.575,1.575)$ before confinement is invoked.
On the four noise bands $[.05,.08]$, $[.08,.10]$, $[.10,.12]$ and
$[.12,\Lambda]$, the finite central amplitude ceilings are
$.375,.325,.275,.25$, and the left ceilings are $.40,.35,.30,.275$.
Each ceiling is intersected with its row's independently acquired
global amplitude. These finite profiles use the same proved
confinement gates with their own parameters; the uniform-tail profiles
are not substituted for them.

The accepted acquisition cover has 184 closed cells: 92 exclusions,
45 left rows covering $[.29,.38]$ and 47 central rows covering
$[.38,.474]$. Thus the moment and probability complements are excluded
before the comparison. Every row retains its own resource and both
threshold signs. The central positive-noise cover uses TINY on $(0,.05]$;
above .05 the whole-count phase radii are .09,.14,.19,.29,.5 on bands
separated by .06,.08,.10,.12, with the complementary nearest-count
phases through .12. The left cover uses TINY on $(0,.05]$, full-phase
MID on $[.05,.08]$, then nearest count outside and whole count inside
phase radii .19 and .27 on $[.08,.10]$ and $[.10,.12]$; above .12 the
whole-count test covers every phase. Each band is intersected with its
actual acquired noise ceiling. The 78 central and 74 left signed
whole-count rows, and the nearest-count covers, include every seam.

The signed cubic multipliers, noise-dependent inward bound, localized
binomial mass, $4553/1000$ harmonic inequality and signed Gaussian
correction proved above give the strict comparison on these entire
closed cells. The two complete LOW families retain their inverse-index
sign gates before any floor substitution. Zero noise is Bernoulli.
The strict local comparisons and Schulz's bound contradict $C_n>\CE$.
This proposition has the stated upper index $48$; it is separate from the
floor-$40$ comparison used only at actual indices $40,41,42$.
\end{proof}

\begin{proposition}\label{j:local:uniform}
There is no attained positive closed-event ratio maximizer satisfying
\begin{gather*}
 n\ge47,\qquad 0\le n(C_n-\CE)\le7/10,\qquad
 n(C_{n-1}-C_n)_+\le19/100,\\
 C_n<4247/10000,\qquad C_{n-1},C_{n-2}<4253/10000.
\end{gather*}
\end{proposition}
\begin{proof}
The support, acquisition and confinement implications proved above and
their accepted closed numerical covers place the same law in one of the
45 central or 38 left rows. Their six local tests cover every positive
noise, both signs of $H$ and every phase, and give
$\rho_B(C_n-C_B)<-.001\lambda$. Schulz's theorem gives $C_B<\CE$.
For zero noise the law is exactly Bernoulli, and the same finite-binomial
theorem gives a strict bound at every threshold. Both cases contradict
$C_n\ge\CE$.
The initial score interval is $[.4095,.4247]$, the native barrier is zero,
and the resource identity uses $e=C_n-\CE$ only as a nonnegative summand.
Strict numerical reserves are independent of $e>0$. The stored value $98$
in initialization occurs only in $.4095+K/98$ with $K=0$; the support's
last band and all later monotonicity gates have no upper index limit.
\end{proof}
\begin{corollary}\label{j:local:equality}
Once $C_j\le\CE$ has been proved at every index, $C_n<\CE$ for every
$n\ge47$, in particular for every $n\ge700$.
\end{corollary}
\begin{proof}
If $C_n=\CE$, ratio attainment supplies the positive closed-event maximum.
The predecessor loss and scaled excess are zero, and both deleted bounds
are at most $\CE$. All hypotheses of the proposition hold. Attainment is
essential here: strictness for each individual law alone would not rule
out a nonattained supremum.
\end{proof}
